% Options for packages loaded elsewhere
\PassOptionsToPackage{unicode}{hyperref}
\PassOptionsToPackage{hyphens}{url}
\documentclass[
  nobib]{tufte-handout}
\usepackage{xcolor}
\usepackage{amsmath,amssymb}
\setcounter{secnumdepth}{-\maxdimen} % remove section numbering
\usepackage{iftex}
\ifPDFTeX
  \usepackage[T1]{fontenc}
  \usepackage[utf8]{inputenc}
  \usepackage{textcomp} % provide euro and other symbols
\else % if luatex or xetex
  \usepackage{unicode-math} % this also loads fontspec
  \defaultfontfeatures{Scale=MatchLowercase}
  \defaultfontfeatures[\rmfamily]{Ligatures=TeX,Scale=1}
\fi
\usepackage{lmodern}
\ifPDFTeX\else
  % xetex/luatex font selection
\fi
% Use upquote if available, for straight quotes in verbatim environments
\IfFileExists{upquote.sty}{\usepackage{upquote}}{}
\IfFileExists{microtype.sty}{% use microtype if available
  \usepackage[]{microtype}
  \UseMicrotypeSet[protrusion]{basicmath} % disable protrusion for tt fonts
}{}
\makeatletter
\@ifundefined{KOMAClassName}{% if non-KOMA class
  \IfFileExists{parskip.sty}{%
    \usepackage{parskip}
  }{% else
    \setlength{\parindent}{0pt}
    \setlength{\parskip}{6pt plus 2pt minus 1pt}}
}{% if KOMA class
  \KOMAoptions{parskip=half}}
\makeatother
\usepackage{color}
\usepackage{fancyvrb}
\newcommand{\VerbBar}{|}
\newcommand{\VERB}{\Verb[commandchars=\\\{\}]}
\DefineVerbatimEnvironment{Highlighting}{Verbatim}{commandchars=\\\{\}}
% Add ',fontsize=\small' for more characters per line
\newenvironment{Shaded}{}{}
\newcommand{\AlertTok}[1]{\textcolor[rgb]{1.00,0.00,0.00}{\textbf{#1}}}
\newcommand{\AnnotationTok}[1]{\textcolor[rgb]{0.38,0.63,0.69}{\textbf{\textit{#1}}}}
\newcommand{\AttributeTok}[1]{\textcolor[rgb]{0.49,0.56,0.16}{#1}}
\newcommand{\BaseNTok}[1]{\textcolor[rgb]{0.25,0.63,0.44}{#1}}
\newcommand{\BuiltInTok}[1]{\textcolor[rgb]{0.00,0.50,0.00}{#1}}
\newcommand{\CharTok}[1]{\textcolor[rgb]{0.25,0.44,0.63}{#1}}
\newcommand{\CommentTok}[1]{\textcolor[rgb]{0.38,0.63,0.69}{\textit{#1}}}
\newcommand{\CommentVarTok}[1]{\textcolor[rgb]{0.38,0.63,0.69}{\textbf{\textit{#1}}}}
\newcommand{\ConstantTok}[1]{\textcolor[rgb]{0.53,0.00,0.00}{#1}}
\newcommand{\ControlFlowTok}[1]{\textcolor[rgb]{0.00,0.44,0.13}{\textbf{#1}}}
\newcommand{\DataTypeTok}[1]{\textcolor[rgb]{0.56,0.13,0.00}{#1}}
\newcommand{\DecValTok}[1]{\textcolor[rgb]{0.25,0.63,0.44}{#1}}
\newcommand{\DocumentationTok}[1]{\textcolor[rgb]{0.73,0.13,0.13}{\textit{#1}}}
\newcommand{\ErrorTok}[1]{\textcolor[rgb]{1.00,0.00,0.00}{\textbf{#1}}}
\newcommand{\ExtensionTok}[1]{#1}
\newcommand{\FloatTok}[1]{\textcolor[rgb]{0.25,0.63,0.44}{#1}}
\newcommand{\FunctionTok}[1]{\textcolor[rgb]{0.02,0.16,0.49}{#1}}
\newcommand{\ImportTok}[1]{\textcolor[rgb]{0.00,0.50,0.00}{\textbf{#1}}}
\newcommand{\InformationTok}[1]{\textcolor[rgb]{0.38,0.63,0.69}{\textbf{\textit{#1}}}}
\newcommand{\KeywordTok}[1]{\textcolor[rgb]{0.00,0.44,0.13}{\textbf{#1}}}
\newcommand{\NormalTok}[1]{#1}
\newcommand{\OperatorTok}[1]{\textcolor[rgb]{0.40,0.40,0.40}{#1}}
\newcommand{\OtherTok}[1]{\textcolor[rgb]{0.00,0.44,0.13}{#1}}
\newcommand{\PreprocessorTok}[1]{\textcolor[rgb]{0.74,0.48,0.00}{#1}}
\newcommand{\RegionMarkerTok}[1]{#1}
\newcommand{\SpecialCharTok}[1]{\textcolor[rgb]{0.25,0.44,0.63}{#1}}
\newcommand{\SpecialStringTok}[1]{\textcolor[rgb]{0.73,0.40,0.53}{#1}}
\newcommand{\StringTok}[1]{\textcolor[rgb]{0.25,0.44,0.63}{#1}}
\newcommand{\VariableTok}[1]{\textcolor[rgb]{0.10,0.09,0.49}{#1}}
\newcommand{\VerbatimStringTok}[1]{\textcolor[rgb]{0.25,0.44,0.63}{#1}}
\newcommand{\WarningTok}[1]{\textcolor[rgb]{0.38,0.63,0.69}{\textbf{\textit{#1}}}}
\usepackage{longtable,booktabs,array}
\newcounter{none} % for unnumbered tables
\usepackage{calc} % for calculating minipage widths
% Correct order of tables after \paragraph or \subparagraph
\usepackage{etoolbox}
\makeatletter
\patchcmd\longtable{\par}{\if@noskipsec\mbox{}\fi\par}{}{}
\makeatother
% Allow footnotes in longtable head/foot
\IfFileExists{footnotehyper.sty}{\usepackage{footnotehyper}}{\usepackage{footnote}}
\makesavenoteenv{longtable}
\setlength{\emergencystretch}{3em} % prevent overfull lines
\providecommand{\tightlist}{%
  \setlength{\itemsep}{0pt}\setlength{\parskip}{0pt}}
% =====================================================================
% axona-doc-preamble.tex — shared LaTeX preamble for the programmer-guide
% quartet (Quick-Start, Programmer Guide, API Reference, Services Guide).
%
% These docs are markdown-native; the PDF is built by converting the .md to
% a tufte-handout .tex with pandoc (-H this file) and compiling with tectonic.
% See render.sh. Modeled on explainer/Axona-Explainer.tex (the flat-essay
% original); the tweaks below make tufte-handout work for code/table-heavy
% reference material.
% =====================================================================

\definecolor{rust}{RGB}{185,78,54}
\definecolor{inkmuted}{RGB}{102,102,102}
\usepackage{microtype}

% XeTeX (tectonic) unicode monospace: renders … → and the box-drawing ASCII
% diagrams that the default CM/LM mono cannot. Menlo ships with macOS.
\setmonofont{Menlo}[Scale=0.78]

% tufte-handout ships \subsubsection as an error stub (it wants flat essays).
% A run-in fallback lets nested reference headings (§4.2, §12.3, …) compile.
\renewcommand{\subsubsection}[1]{\medskip\par\noindent{\normalsize\bfseries #1}\par\nobreak\smallskip}

% These docs use no margin notes, so reclaim tufte's wide right margin: widen
% the text block so code blocks and API tables don't overflow it.
\usepackage{geometry}
\geometry{left=1in, textwidth=6.1in, marginparsep=0.2in, marginparwidth=0.9in,
          top=1in, headsep=0.4in}

% Wrap long code lines (cheat-sheet comments run to ~130 chars) with a ↪ marker
% instead of overflowing. Applies to pandoc's fancyvrb Highlighting/Verbatim.
\usepackage{fvextra}
\fvset{breaklines=true, breakanywhere=true}

% A few technical symbols appear in prose / inline code that Latin Modern and
% Menlo don't cover; map them to math equivalents so they render (rather than
% dropping to tofu). Only maps chars actually used — leaves → alone (it renders).
\usepackage{newunicodechar}
\newunicodechar{≈}{\ensuremath{\approx}}
\newunicodechar{≥}{\ensuremath{\geq}}
\newunicodechar{≤}{\ensuremath{\leq}}
\newunicodechar{×}{\ensuremath{\times}}
\newunicodechar{·}{\ensuremath{\cdot}}
\newunicodechar{‖}{\ensuremath{\|}}
\usepackage{bookmark}
\IfFileExists{xurl.sty}{\usepackage{xurl}}{} % add URL line breaks if available
\urlstyle{same}
\hypersetup{
  hidelinks,
  pdfcreator={LaTeX via pandoc}}

\author{}
\date{}

\begin{document}

\section{Axona Programmer Guide}\label{axona-programmer-guide}

A practical guide to building a decentralized pub/sub application on the
Axona protocol. After reading this you should be able to ship a working
chat, feed, inbox, or notification system without reading any other
Axona code.

The guide assumes you already know JavaScript + Node + a browser. It
does not assume any DHT / WebRTC / cryptography background -- concepts
are introduced where they are needed.

\begin{itemize}
\tightlist
\item
  \textbf{Protocol kernel}:
  \href{https://github.com/axona-net/axona-protocol}{@axona/protocol}
  (v4.11.2)
\item
  \textbf{Wire version}: 4.0 (\texttt{WIRE\_VERSION}); kernel version
  4.11.2 (\texttt{KERNEL\_VERSION})
\item
  \textbf{Live network}: the 4.x line runs on \textbf{testnet} ---
  \texttt{wss://testnet.axona.net}. The \textbf{production} bridges
  (\texttt{wss://bridge.axona.net} east,
  \texttt{wss://bridge-west.axona.net} west --- a federated pair) are
  still on the 3.x line and do \textbf{not} interoperate with 4.x (the
  wire major partitions them). Build against testnet to use this
  surface.
\item
  \textbf{Companion docs}:

  \begin{itemize}
  \tightlist
  \item
    \href{Quick-Start-v4.11.2.md}{Quick Start} -- five minutes to a
    working roundtrip; send a newcomer there first.
  \item
    \href{Axona-API-Reference-v4.11.2.md}{API Reference} -- every public
    symbol and its exact signature.
  \end{itemize}
\end{itemize}

\begin{quote}
\textbf{What changed from the v2 line.} v3.0.0 rebuilt identity,
authorship, and addressing as three separate concerns (a breaking
flag-day). If you learned Axona on the v2 guide, the things you must
unlearn are: the single \texttt{deriveIdentity}, the
\texttt{publisher}/\texttt{publishId} arguments, \texttt{sign:\ false},
bare string topics, and the ``region-keyed / publisher-keyed / public''
topic \emph{modes}. All gone. The replacements are two identity
factories (\texttt{createNodeIdentity}, \texttt{createAuthorIdentity}),
descriptor topics (\texttt{\{\ region?,\ owner?,\ name,\ write?\ \}}),
and per-publish \texttt{signWith}.

\textbf{What changed in the 4.x line.} The application API above is
\textbf{unchanged} --- the 4.x work is a routing/reliability rewrite
underneath it: routing-only axonic-tree pub/sub, K-closest
\textbf{cohort} replication (a killed message can't resurface and a
publish survives the root churning out), a \textbf{cold-publish burst}
so a freshly-joined node's first publish isn't lost,
\textbf{nearest-replica reads} (\texttt{pull} answers from the closest
replica holding the message), and a rebuilt \textbf{metrics} path. All
transparent to your code; the reliability notes appear in §7 where they
matter.
\end{quote}

\begin{center}\rule{0.5\linewidth}{0.5pt}\end{center}

\subsection{Table of contents}\label{table-of-contents}

\begin{enumerate}
\def\labelenumi{\arabic{enumi}.}
\tightlist
\item
  \hyperref[1-what-is-axona]{What is Axona?}
\item
  \hyperref[2-setup]{Setup}
\item
  \hyperref[3-mental-model]{Mental model}
\item
  \hyperref[4-node-identity-the-connection]{Node identity (the
  connection)}
\item
  \hyperref[5-author-identity-and-signing-the-authorship]{Author
  identity and signing (the authorship)}
\item
  \hyperref[6-topics-and-addressing]{Topics and addressing}
\item
  \hyperref[7-publish-and-subscribe]{Publish and subscribe}
\item
  \hyperref[8-direct-messaging]{Direct messaging}
\item
  \hyperref[9-mesh-introspection]{Mesh introspection}
\item
  \hyperref[10-persistence]{Persistence}
\item
  \hyperref[11-the-bridge-docker-production]{The bridge (Docker
  production)}
\item
  \hyperref[12-worked-example-a-regional-chat-app]{Worked example: a
  regional chat app}
\item
  \hyperref[13-common-pitfalls]{Common pitfalls}
\item
  \hyperref[14-production-checklist]{Production checklist}
\item
  \hyperref[15-cheat-sheet]{Cheat sheet}
\end{enumerate}

\begin{center}\rule{0.5\linewidth}{0.5pt}\end{center}

\subsection{1. What is Axona?}\label{what-is-axona}

Axona is a peer-to-peer mesh protocol where every participant is an
equal node -- there is no central server. Three things make it useful:

\begin{enumerate}
\def\labelenumi{\arabic{enumi}.}
\tightlist
\item
  \textbf{Geographic routing.} Every node's ID and every topic's ID
  begins with a one-byte region prefix (a Google S2 cell). Lookups and
  message placement are XOR-routed in a way that strongly favors
  same-region neighbors, so a topic pinned to a region naturally
  clusters on nodes in that region.
\item
  \textbf{Signed authorship.} Every published message is signed by an
  \emph{author key} (Ed25519). A reader can prove who authored a
  message, that it was not tampered with, and -- if the author chose to
  be durable -- recognize the same author later.
\item
  \textbf{Replicated cached delivery.} Each topic has a set of
  \emph{axons} -- the K closest peers to the topic ID -- that hold a
  bounded ring of recent messages. Late subscribers can replay history;
  live subscribers tail in real time.
\end{enumerate}

The protocol kernel (\texttt{@axona/protocol}) is environment-neutral
pure JS. It runs in browsers, Node, and in-process simulation. Deployed
components ride on top of it:

{\def\LTcaptype{none} % do not increment counter
\begin{longtable}[]{@{}
  >{\raggedright\arraybackslash}p{(\linewidth - 4\tabcolsep) * \real{0.3333}}
  >{\raggedright\arraybackslash}p{(\linewidth - 4\tabcolsep) * \real{0.3333}}
  >{\raggedright\arraybackslash}p{(\linewidth - 4\tabcolsep) * \real{0.3333}}@{}}
\toprule\noalign{}
\begin{minipage}[b]{\linewidth}\raggedright
Component
\end{minipage} & \begin{minipage}[b]{\linewidth}\raggedright
What it is
\end{minipage} & \begin{minipage}[b]{\linewidth}\raggedright
Where it runs
\end{minipage} \\
\midrule\noalign{}
\endhead
\bottomrule\noalign{}
\endlastfoot
\texttt{axona-peer} & Browser SDK + reference client & Static HTML/JS,
any web host \\
\texttt{axona-bridge} & WebRTC signaling broker + universal-hub peer & A
Docker stack (bridge + Caddy + coturn) behind auto-TLS \\
\texttt{axona-relay} & Headless hosting node (stores + serves topics) &
Node \\
\texttt{dht-sim} & Visualization + analysis simulator & Node, locally \\
\end{longtable}
}

The bridge is a signaling broker for WebRTC and a universal-fallback
peer in the mesh -- it is \textbf{not} a centralized message queue. Most
messages flow directly between browsers over WebRTC data channels; the
bridge is a WebSocket fallback when NAT prevents a direct connection,
and a hub that helps mesh discovery on cold start. The protocol does not
depend on the bridge for correctness.

\subsubsection{The three separated concerns (read this once and it all
fits)}\label{the-three-separated-concerns-read-this-once-and-it-all-fits}

The single most important idea in Axona v3 is that \textbf{connection,
authorship, and addressing are three different things, each with its own
primitive, and no one of them does another's job:}

{\def\LTcaptype{none} % do not increment counter
\begin{longtable}[]{@{}
  >{\raggedright\arraybackslash}p{(\linewidth - 6\tabcolsep) * \real{0.2500}}
  >{\raggedright\arraybackslash}p{(\linewidth - 6\tabcolsep) * \real{0.2500}}
  >{\raggedright\arraybackslash}p{(\linewidth - 6\tabcolsep) * \real{0.2500}}
  >{\raggedright\arraybackslash}p{(\linewidth - 6\tabcolsep) * \real{0.2500}}@{}}
\toprule\noalign{}
\begin{minipage}[b]{\linewidth}\raggedright
Concern
\end{minipage} & \begin{minipage}[b]{\linewidth}\raggedright
The question it answers
\end{minipage} & \begin{minipage}[b]{\linewidth}\raggedright
Primitive
\end{minipage} & \begin{minipage}[b]{\linewidth}\raggedright
Lifetime
\end{minipage} \\
\midrule\noalign{}
\endhead
\bottomrule\noalign{}
\endlastfoot
\textbf{Connection} & ``Which node is this, and is it really it on the
wire?'' & \textbf{Node Identity} -\textgreater{} \textbf{Node ID} &
ephemeral (fresh each run) \\
\textbf{Authorship} & ``Who is accountable for this message, provably?''
& \textbf{Author Identity} -\textgreater{} \textbf{Author ID} & your
choice: per-run or durable \\
\textbf{Addressing} & ``Where does this topic live? Is this a
duplicate?'' & \textbf{Topic ID} / \textbf{Message ID} & deterministic,
identity-independent \\
\end{longtable}
}

Two consequences are the heart of the model:

\begin{itemize}
\tightlist
\item
  \textbf{Connection is not authorship.} The connection key is thrown
  away and re-minted each run (so being online is not trackable). The
  author key is the thing you keep (so your authorship is recognizable).
  Either can change without disturbing the other.
\item
  \textbf{Addressing is not identity.} Where a message lives and whether
  it is a duplicate come from the topic and the content, not from who
  you are. Routing and dedup work for anonymous senders too.
\end{itemize}

Sections 4, 5, and 6 are one per concern. Internalize the table above
and the rest is detail.

\subsubsection{Address space}\label{address-space}

Every Node ID and every Topic ID is \textbf{264 bits = 66 lowercase hex
characters}. The layout is the same for both: a one-byte region prefix
followed by a 256-bit hash.

\begin{verbatim}
   +--------------------------+---------------------------------------+
   |  region byte             |  256-bit hash                         |
   |  8 bits (top 2 hex)      |  64 hex chars                         |
   +--------------------------+---------------------------------------+
\end{verbatim}

\begin{itemize}
\tightlist
\item
  \textbf{Node ID} =
  \texttt{regionByte\ \textbar{}\textbar{}\ SHA-256(pubkey)} -- the
  region is the S2 cell of the node's location; the hash is of its
  connection public key.
\item
  \textbf{Topic ID} =
  \texttt{regionByte\ \textbar{}\textbar{}\ SHA-256(canonical(\{\ owner,\ name,\ write\ \}))}
  -- see section 6.
\end{itemize}

Because node IDs and topic IDs share the same space, the XOR distance
between a node and a topic is meaningful, and the router uses it to keep
locality.

\begin{center}\rule{0.5\linewidth}{0.5pt}\end{center}

\subsection{2. Setup}\label{setup}

\subsubsection{2.1 Install the kernel}\label{install-the-kernel}

\begin{verbatim}
mkdir my-axona-app && cd my-axona-app
npm init -y
npm install @axona/protocol@github:axona-net/axona-protocol#v4.11.2
\end{verbatim}

You now have \texttt{node\_modules/@axona/protocol/src/} with the full
kernel.

\subsubsection{2.2 The shape of a peer}\label{the-shape-of-a-peer}

A running peer is four objects:

\begin{Shaded}
\begin{Highlighting}[]
\ImportTok{import}\NormalTok{ \{}
\NormalTok{  AxonaPeer}\OperatorTok{,}\NormalTok{ AxonaDomain}\OperatorTok{,}\NormalTok{ NeuronNode}\OperatorTok{,}
\NormalTok{  createNodeIdentity}\OperatorTok{,}\NormalTok{ createAuthorIdentity}\OperatorTok{,}
\NormalTok{  KERNEL\_VERSION}\OperatorTok{,}
\NormalTok{\} }\ImportTok{from} \StringTok{\textquotesingle{}@axona/protocol\textquotesingle{}}\OperatorTok{;}
\ImportTok{import}\NormalTok{ \{ webTransport \} }\ImportTok{from} \StringTok{\textquotesingle{}@axona/protocol/transport/web/index.js\textquotesingle{}}\OperatorTok{;}

\CommentTok{// 1. The CONNECTION identity (needs a location {-}\textgreater{} yields a Node ID).}
\KeywordTok{const}\NormalTok{ nodeIdentity }\OperatorTok{=} \ControlFlowTok{await} \FunctionTok{createNodeIdentity}\NormalTok{(\{ }\DataTypeTok{lat}\OperatorTok{:} \FloatTok{38.0}\OperatorTok{,} \DataTypeTok{lng}\OperatorTok{:} \OperatorTok{{-}}\FloatTok{77.0}\NormalTok{ \})}\OperatorTok{;}

\CommentTok{// 2. A transport. In the browser, webTransport speaks WebRTC + a WS bridge.}
\KeywordTok{const}\NormalTok{ transport }\OperatorTok{=} \FunctionTok{webTransport}\NormalTok{(\{}
  \DataTypeTok{bridgeUrl}\OperatorTok{:} \StringTok{\textquotesingle{}wss://testnet.axona.net\textquotesingle{}}\OperatorTok{,}  \CommentTok{// the 4.x line runs on testnet (prod is 3.x)}
  \DataTypeTok{identity}\OperatorTok{:}\NormalTok{  nodeIdentity}\OperatorTok{,}               \CommentTok{// the factory option is named \textasciigrave{}identity:\textasciigrave{}}
\NormalTok{\})}\OperatorTok{;}

\CommentTok{// 3. The local routing state.}
\KeywordTok{const}\NormalTok{ node }\OperatorTok{=} \KeywordTok{new} \FunctionTok{NeuronNode}\NormalTok{(\{}
  \DataTypeTok{id}\OperatorTok{:}  \BuiltInTok{BigInt}\NormalTok{(}\StringTok{\textquotesingle{}0x\textquotesingle{}} \OperatorTok{+}\NormalTok{ nodeIdentity}\OperatorTok{.}\AttributeTok{id}\NormalTok{)}\OperatorTok{,}
  \DataTypeTok{lat}\OperatorTok{:} \FloatTok{38.0}\OperatorTok{,} \DataTypeTok{lng}\OperatorTok{:} \OperatorTok{{-}}\FloatTok{77.0}\OperatorTok{,}
\NormalTok{\})}\OperatorTok{;}
\NormalTok{node}\OperatorTok{.}\AttributeTok{transport} \OperatorTok{=}\NormalTok{ transport}\OperatorTok{;}

\CommentTok{// 4. The peer ties them together. Note: NO author key here {-}{-} authorship}
\CommentTok{//    is per{-}publish, never a constructor argument (see section 5).}
\KeywordTok{const}\NormalTok{ peer }\OperatorTok{=} \KeywordTok{new} \FunctionTok{AxonaPeer}\NormalTok{(\{}
  \DataTypeTok{domain}\OperatorTok{:}       \KeywordTok{new} \FunctionTok{AxonaDomain}\NormalTok{(\{ }\DataTypeTok{k}\OperatorTok{:} \DecValTok{20}\NormalTok{ \})}\OperatorTok{,}
\NormalTok{  node}\OperatorTok{,}
\NormalTok{  nodeIdentity}\OperatorTok{,}
\NormalTok{  transport}\OperatorTok{,}
\NormalTok{\})}\OperatorTok{;}

\ControlFlowTok{await}\NormalTok{ transport}\OperatorTok{.}\FunctionTok{start}\NormalTok{(nodeIdentity}\OperatorTok{.}\AttributeTok{id}\NormalTok{)}\OperatorTok{;}
\ControlFlowTok{await}\NormalTok{ peer}\OperatorTok{.}\FunctionTok{start}\NormalTok{()}\OperatorTok{;}
\BuiltInTok{console}\OperatorTok{.}\FunctionTok{log}\NormalTok{(}\StringTok{\textquotesingle{}connected; kernel\textquotesingle{}}\OperatorTok{,}\NormalTok{ KERNEL\_VERSION)}\OperatorTok{;}
\end{Highlighting}
\end{Shaded}

In Node you would use a node transport instead of \texttt{webTransport},
but the peer-construction shape is identical. The fastest way to see all
of this working end to end is to read \texttt{apps/axona-minimal/} in
the kernel repo -- a complete \textasciitilde70-line app that section 12
is based on.

\subsubsection{2.3 Or just open the deployed
peer}\label{or-just-open-the-deployed-peer}

To watch Axona work without writing code, open \url{https://axona.net}
in two browser tabs. Each tab is an independent peer. Type the same
topic in both, publish from one, and the message arrives in the other.

\begin{center}\rule{0.5\linewidth}{0.5pt}\end{center}

\subsection{3. Mental model}\label{mental-model}

\subsubsection{3.1 Peers and the mesh}\label{peers-and-the-mesh}

Every participant is a \textbf{peer} -- a long-lived process with:

\begin{itemize}
\tightlist
\item
  \textbf{A node identity}: a 264-bit Node ID and an Ed25519 connection
  keypair.
\item
  \textbf{A synaptome}: an adaptive routing table of other peers,
  weighted by recency, latency, and usefulness.
\item
  \textbf{A transport}: a way to talk to other peers. In the browser
  this is a composite of a WebRTC mesh (direct data channels) and a
  WebSocket fallback through the bridge.
\end{itemize}

Peers form a \textbf{mesh} -- a graph of bidirectional channels. The
mesh is \emph{not} a global broadcast topology. Each peer directly knows
only a handful of others; everything else is reached by routed
forwarding.

The mesh is \textbf{self-healing}, and you don't manage it. A peer
continuously keeps its routing table \emph{complete} in two ways: it
holds its few XOR-nearest ``successor'' peers (so routed delivery's last
hop always lands on the right node) and a spread of longer-range links
(so any target is reachable). When a neighbour drops, the replacement is
almost always a peer it already knows of -- so repair is a cheap local
step, not a global search. As an app author there is nothing to
configure: subscribe, publish, and let the mesh maintain itself
underneath you.

\subsubsection{3.2 Topics, axons, and
replication}\label{topics-axons-and-replication}

A \textbf{topic} is a small descriptor that hashes to a 264-bit Topic ID
(section 6). When you \texttt{peer.pub(...)}, the protocol:

\begin{enumerate}
\def\labelenumi{\arabic{enumi}.}
\tightlist
\item
  Resolves the descriptor to a Topic ID.
\item
  Finds the \textbf{K closest peers} to that Topic ID in XOR space (K
  defaults to 5).
\item
  Sends a publish frame to each. Each appends the message to its local
  replay cache for the topic and forwards it to any subscribers it
  tracks as children.
\end{enumerate}

Those K peers are the topic's \textbf{axons}. They are not a privileged
server set -- any peer that happens to land in the K-closest set for a
Topic ID acts as an axon for it, and the set rotates as nodes join and
leave.

A \texttt{peer.sub(...)} registers the caller as a child of the topic's
axons; from then on, publishes fan out to it. If it asks for
\texttt{\{\ since:\ \textquotesingle{}all\textquotesingle{}\ \}}, the
axon also replays its cache.

\subsubsection{3.3 Geographic clustering}\label{geographic-clustering}

The region byte at the top of every ID means a topic pinned to
\texttt{useast} gets a Topic ID beginning \texttt{0x89}, and
\texttt{useast} nodes -- XOR-close to that prefix -- are naturally
selected as its axons, where its subscribers are also likely to live.
That minimizes cross-region hops. You pin a topic to a region by naming
\texttt{region} in the descriptor (section 6).

\subsubsection{3.4 The signed envelope}\label{the-signed-envelope}

Every publish is wrapped in an envelope and delivered to subscribers as:

\begin{Shaded}
\begin{Highlighting}[]
\NormalTok{\{}
  \DataTypeTok{msgId}\OperatorTok{:}        \StringTok{\textquotesingle{}\textless{}64{-}char hex\textgreater{}\textquotesingle{}}\OperatorTok{,}   \CommentTok{// SHA{-}256(authorId + message); dedup key}
  \DataTypeTok{ts}\OperatorTok{:}           \DecValTok{1779306455550}\OperatorTok{,}     \CommentTok{// ms epoch}
  \DataTypeTok{topic}\OperatorTok{:}\NormalTok{        \{ region}\OperatorTok{,}\NormalTok{ owner}\OperatorTok{,}\NormalTok{ name}\OperatorTok{,}\NormalTok{ write \}}\OperatorTok{,}  \CommentTok{// the SIGNED descriptor (object, not a string)}
  \DataTypeTok{message}\OperatorTok{:}      \OperatorTok{\textless{}}\NormalTok{any }\BuiltInTok{JSON}\NormalTok{ value}\OperatorTok{\textgreater{},}  \CommentTok{// your opaque payload}
  \DataTypeTok{signerPubkey}\OperatorTok{:} \StringTok{\textquotesingle{}\textless{}64{-}char hex\textgreater{}\textquotesingle{}}\OperatorTok{,}   \CommentTok{// the Author ID; absent for anonymous}
  \DataTypeTok{signature}\OperatorTok{:}    \StringTok{\textquotesingle{}ed25519:\textless{}...\textgreater{}\textquotesingle{}}\OperatorTok{,}   \CommentTok{// absent for anonymous}
\NormalTok{\}}
\end{Highlighting}
\end{Shaded}

\texttt{message} is opaque bytes to the protocol -- string, number,
object, array, anything JSON-serializable. The kernel JSON-encodes the
whole envelope at the wire layer.

\subsubsection{3.5 What the protocol secures (and what it does
not)}\label{what-the-protocol-secures-and-what-it-does-not}

Axona is \textbf{self-authenticating}: every guarantee is enforced by
cryptography the peers carry themselves -- no certificate authority, no
central trust server.

\begin{itemize}
\tightlist
\item
  \textbf{Identity is provable, not asserted.} A Node ID's bottom 256
  bits are \texttt{SHA-256(pubkey)}, and every connection runs a
  handshake that proves the peer holds the key and binds the proof to
  \emph{that live connection} so it cannot be replayed onto another
  link.
\item
  \textbf{Authorship is provable and self-contained.} Every accountable
  message carries its Author ID plus a signature over the message;
  receivers verify it and recompute the Message ID. A forwarder can
  never make one author's content look like another's.
\item
  \textbf{Write policy is enforced at the storing node.} An owned topic
  (section 6) accepts a publish only if the signer is the owner
  (\texttt{WRITE\_POLICY\_VIOLATION} otherwise) -- checked at ingress,
  before fan-out.
\item
  \textbf{You can only act for yourself.} A peer subscribes only
  \emph{itself}; nodes host only topics they are routing-responsible
  for; a \texttt{kill} is accepted only if signed by the original
  author.
\item
  \textbf{Location is never disclosed by the protocol.} No message ever
  carries a region, coordinate, or Node ID. The region byte lives only
  inside the ephemeral Node ID, and it is never derived from the author
  key. (If your app wants to show a sender's region it must put that in
  the payload itself -- a deliberate app choice, not a protocol field.
  The worked example in section 12 does exactly this.)
\end{itemize}

\textbf{What it does not do:} encrypt your payload (\texttt{message} is
opaque -- encrypt inside it end-to-end if you need confidentiality), and
it does not decide \emph{whether you trust} a given author -- that is
your app's policy (allowlist, web of trust, trust-on-first-use, or
``accept everyone''). See the \href{../SECURITY-CHANGELOG.md}{security
changelog} for the versioned record.

\begin{center}\rule{0.5\linewidth}{0.5pt}\end{center}

\subsection{4. Node identity (the
connection)}\label{node-identity-the-connection}

\subsubsection{4.1 Creating one}\label{creating-one}

\begin{Shaded}
\begin{Highlighting}[]
\ImportTok{import}\NormalTok{ \{ createNodeIdentity \} }\ImportTok{from} \StringTok{\textquotesingle{}@axona/protocol\textquotesingle{}}\OperatorTok{;}

\KeywordTok{const}\NormalTok{ node }\OperatorTok{=} \ControlFlowTok{await} \FunctionTok{createNodeIdentity}\NormalTok{(\{ }\DataTypeTok{lat}\OperatorTok{:} \FloatTok{38.0}\OperatorTok{,} \DataTypeTok{lng}\OperatorTok{:} \OperatorTok{{-}}\FloatTok{77.0}\NormalTok{ \})}\OperatorTok{;}
\CommentTok{// node = \{}
\CommentTok{//   id:         \textquotesingle{}df27c92b...\textquotesingle{} (66{-}char hex Node ID),}
\CommentTok{//   pubkey:     Uint8Array(32),}
\CommentTok{//   pubkeyHex:  \textquotesingle{}\textless{}64 hex\textgreater{}\textquotesingle{},}
\CommentTok{//   privateKey: \textless{}Web Crypto Ed25519 CryptoKey\textgreater{},}
\CommentTok{//   region:     \{ lat: 38.0, lng: {-}77.0 \},}
\CommentTok{//   createdAt:  1779306455550,}
\CommentTok{//   sign, verify,}
\CommentTok{// \}}
\end{Highlighting}
\end{Shaded}

The top byte of \texttt{node.id} is the S2 cell for
\texttt{(lat,\ lng)}; the bottom 256 bits are \texttt{SHA-256(pubkey)}.
This key authenticates the connection handshake, drives DHT routing, and
is what \texttt{subscribe} rides on. \textbf{It never signs a published
message} (key separation -- section 5).

\texttt{createNodeIdentity} mints a fresh keypair every call. By default
the key is extractable so it can be persisted; pass
\texttt{\{\ extractable:\ false\ \}} for an ephemeral browser identity
that XSS cannot exfiltrate.

\subsubsection{4.2 Ephemeral by default -- and that is the
point}\label{ephemeral-by-default-and-that-is-the-point}

A node identity is meant to be \textbf{thrown away and re-minted each
run.} A restarted peer rejoins as a new node with a new Node ID. This is
a privacy property: nobody can build a durable record of ``this node was
online again.''

You \emph{may} persist it for a stable address (\texttt{dumpIdentity} /
\texttt{loadIdentity}, section 10), but you rarely want to -- a stable
Node ID means linkable connections. The thing you keep across sessions
is your \textbf{author} key (section 5), not your node key.

\subsubsection{4.3 Choosing a location}\label{choosing-a-location}

The location determines your region byte, which is a routing
neighborhood. Three strategies:

{\def\LTcaptype{none} % do not increment counter
\begin{longtable}[]{@{}
  >{\raggedright\arraybackslash}p{(\linewidth - 2\tabcolsep) * \real{0.5000}}
  >{\raggedright\arraybackslash}p{(\linewidth - 2\tabcolsep) * \real{0.5000}}@{}}
\toprule\noalign{}
\begin{minipage}[b]{\linewidth}\raggedright
Strategy
\end{minipage} & \begin{minipage}[b]{\linewidth}\raggedright
When
\end{minipage} \\
\midrule\noalign{}
\endhead
\bottomrule\noalign{}
\endlastfoot
Hardcode \texttt{\{\ lat,\ lng\ \}} & A server-side peer (relay, daemon,
scraper) whose region is known. \\
Geolocate the user & \texttt{navigator.geolocation.getCurrentPosition},
falling back to a default region if denied. \\
Let the user pick & A region dropdown; use \texttt{regionCenter(name)}
from the region module to turn a choice into
\texttt{\{\ lat,\ lng\ \}}. \\
\end{longtable}
}

The region module (\texttt{@axona/protocol}) ships the canonical region
table and converters:

\begin{Shaded}
\begin{Highlighting}[]
\ImportTok{import}\NormalTok{ \{ regionName}\OperatorTok{,}\NormalTok{ regionCode}\OperatorTok{,}\NormalTok{ resolveRegion}\OperatorTok{,}\NormalTok{ regionCenter \} }\ImportTok{from} \StringTok{\textquotesingle{}@axona/protocol\textquotesingle{}}\OperatorTok{;}

\FunctionTok{regionName}\NormalTok{(}\BaseNTok{0x89}\NormalTok{)}\OperatorTok{;}        \CommentTok{// {-}\textgreater{} \textquotesingle{}useast\textquotesingle{}}
\FunctionTok{regionCode}\NormalTok{(}\StringTok{\textquotesingle{}useast\textquotesingle{}}\NormalTok{)}\OperatorTok{;}    \CommentTok{// {-}\textgreater{} 0x89 (137)}
\FunctionTok{resolveRegion}\NormalTok{(}\StringTok{\textquotesingle{}0x89\textquotesingle{}}\NormalTok{)}\OperatorTok{;}   \CommentTok{// {-}\textgreater{} 137   (accepts a name, \textquotesingle{}0x89\textquotesingle{}, \textquotesingle{}137\textquotesingle{}, or 137)}
\FunctionTok{regionCenter}\NormalTok{(}\StringTok{\textquotesingle{}useast\textquotesingle{}}\NormalTok{)}\OperatorTok{;}  \CommentTok{// {-}\textgreater{} \{ lat, lng \} {-}{-} feed straight into createNodeIdentity}
\end{Highlighting}
\end{Shaded}

Your node region is independent of which regions' topics you address --
you can publish to a topic pinned to any region regardless of where your
node lives (section 6).

\begin{center}\rule{0.5\linewidth}{0.5pt}\end{center}

\subsection{5. Author identity and signing (the
authorship)}\label{author-identity-and-signing-the-authorship}

\subsubsection{5.1 The author key has no location and no Node
ID}\label{the-author-key-has-no-location-and-no-node-id}

\begin{Shaded}
\begin{Highlighting}[]
\ImportTok{import}\NormalTok{ \{ createAuthorIdentity \} }\ImportTok{from} \StringTok{\textquotesingle{}@axona/protocol\textquotesingle{}}\OperatorTok{;}

\KeywordTok{const}\NormalTok{ me }\OperatorTok{=} \ControlFlowTok{await} \FunctionTok{createAuthorIdentity}\NormalTok{()}\OperatorTok{;}          \CommentTok{// ephemeral author}
\CommentTok{// me = \{ kind: \textquotesingle{}author\textquotesingle{}, authorId: \textquotesingle{}\textless{}64 hex\textgreater{}\textquotesingle{}, pubkey, pubkeyHex, privateKey, sign, verify \}}
\end{Highlighting}
\end{Shaded}

\texttt{me.authorId} is your public \textbf{Author ID} -- the value that
appears on every message you sign (it is the \texttt{signerPubkey} field
on the wire). An author identity is a keypair and nothing else:
\textbf{no location, no region, no Node ID.} Authorship is not a place.

\subsubsection{5.2 Durability is the author's choice -- the only
persistence decision most apps
make}\label{durability-is-the-authors-choice-the-only-persistence-decision-most-apps-make}

\begin{Shaded}
\begin{Highlighting}[]
\ControlFlowTok{await} \FunctionTok{createAuthorIdentity}\NormalTok{()}\OperatorTok{;}                          \CommentTok{// ephemeral: one{-}shot, unlinkable}
\ControlFlowTok{await} \FunctionTok{createAuthorIdentity}\NormalTok{(\{ }\DataTypeTok{persistAs}\OperatorTok{:} \StringTok{\textquotesingle{}me\textquotesingle{}}\NormalTok{ \})}\OperatorTok{;}       \CommentTok{// durable: load{-}or{-}create + save to localStorage}
\ControlFlowTok{await} \FunctionTok{createAuthorIdentity}\NormalTok{(\{ }\DataTypeTok{persistAs}\OperatorTok{:} \StringTok{\textquotesingle{}me\textquotesingle{}}\OperatorTok{,}\NormalTok{ store \})}\OperatorTok{;} \CommentTok{// durable: custom \{ get, set \} store}
\end{Highlighting}
\end{Shaded}

\begin{itemize}
\tightlist
\item
  \textbf{Ephemeral} -- minted fresh, unlinkable, gone when the tab
  closes. Good for genuinely throwaway posting.
\item
  \textbf{Durable} (\texttt{persistAs}) -- on first call it mints and
  saves; on later calls it loads the saved key. The same Author ID
  survives reloads, so readers recognize you across sessions and you can
  \texttt{kill} your own content later. \texttt{persistAs} is a
  \emph{local storage label}, not a network name.
\end{itemize}

In the browser \texttt{persistAs} uses \texttt{localStorage} by default;
pass \texttt{store} (any \texttt{\{\ get,\ set\ \}}) to control where it
lands. A persisted key is forced extractable so it can be serialized.

\subsubsection{5.3 An app may hold several authors -- or
none}\label{an-app-may-hold-several-authors-or-none}

A peer holds \textbf{no default author.} Authorship is supplied at the
one moment it is used -- signing a publish -- via \texttt{signWith}.
This is deliberate:

\begin{itemize}
\tightlist
\item
  A subscribe-only peer needs no author at all.
\item
  A peer can run several authors (personas) at once, choosing which
  signs each message.
\end{itemize}

\begin{Shaded}
\begin{Highlighting}[]
\KeywordTok{const}\NormalTok{ alice }\OperatorTok{=} \ControlFlowTok{await} \FunctionTok{createAuthorIdentity}\NormalTok{(\{ }\DataTypeTok{persistAs}\OperatorTok{:} \StringTok{\textquotesingle{}alice\textquotesingle{}}\NormalTok{ \})}\OperatorTok{;}
\KeywordTok{const}\NormalTok{ bot   }\OperatorTok{=} \ControlFlowTok{await} \FunctionTok{createAuthorIdentity}\NormalTok{(\{ }\DataTypeTok{persistAs}\OperatorTok{:} \StringTok{\textquotesingle{}bot\textquotesingle{}}\NormalTok{ \})}\OperatorTok{;}

\ControlFlowTok{await}\NormalTok{ peer}\OperatorTok{.}\FunctionTok{pub}\NormalTok{(\{ }\DataTypeTok{name}\OperatorTok{:} \StringTok{\textquotesingle{}lobby\textquotesingle{}}\NormalTok{ \}}\OperatorTok{,} \StringTok{\textquotesingle{}hi from alice\textquotesingle{}}\OperatorTok{,}\NormalTok{ \{ }\DataTypeTok{signWith}\OperatorTok{:}\NormalTok{ alice \})}\OperatorTok{;}
\ControlFlowTok{await}\NormalTok{ peer}\OperatorTok{.}\FunctionTok{pub}\NormalTok{(\{ }\DataTypeTok{name}\OperatorTok{:} \StringTok{\textquotesingle{}lobby\textquotesingle{}}\NormalTok{ \}}\OperatorTok{,} \StringTok{\textquotesingle{}beep boop\textquotesingle{}}\OperatorTok{,}\NormalTok{     \{ }\DataTypeTok{signWith}\OperatorTok{:}\NormalTok{ bot \})}\OperatorTok{;}
\end{Highlighting}
\end{Shaded}

\subsubsection{5.4 Every publish names its signer -- there is no
default}\label{every-publish-names-its-signer-there-is-no-default}

This is a hard rule. Each \texttt{peer.pub} (and \texttt{kill}) takes a
\texttt{signWith}:

\begin{Shaded}
\begin{Highlighting}[]
\ImportTok{import}\NormalTok{ \{ ANONYMOUS \} }\ImportTok{from} \StringTok{\textquotesingle{}@axona/protocol\textquotesingle{}}\OperatorTok{;}

\ControlFlowTok{await}\NormalTok{ peer}\OperatorTok{.}\FunctionTok{pub}\NormalTok{(\{ }\DataTypeTok{name}\OperatorTok{:} \StringTok{\textquotesingle{}lobby\textquotesingle{}}\NormalTok{ \}}\OperatorTok{,}\NormalTok{ msg}\OperatorTok{,}\NormalTok{ \{ }\DataTypeTok{signWith}\OperatorTok{:}\NormalTok{ me \})}\OperatorTok{;}         \CommentTok{// signed by \textasciigrave{}me\textasciigrave{}}
\ControlFlowTok{await}\NormalTok{ peer}\OperatorTok{.}\FunctionTok{pub}\NormalTok{(\{ }\DataTypeTok{name}\OperatorTok{:} \StringTok{\textquotesingle{}lobby\textquotesingle{}}\NormalTok{ \}}\OperatorTok{,}\NormalTok{ msg}\OperatorTok{,}\NormalTok{ \{ }\DataTypeTok{signWith}\OperatorTok{:}\NormalTok{ ANONYMOUS \})}\OperatorTok{;}  \CommentTok{// explicitly anonymous (unsigned)}
\ControlFlowTok{await}\NormalTok{ peer}\OperatorTok{.}\FunctionTok{pub}\NormalTok{(\{ }\DataTypeTok{name}\OperatorTok{:} \StringTok{\textquotesingle{}lobby\textquotesingle{}}\NormalTok{ \}}\OperatorTok{,}\NormalTok{ msg)}\OperatorTok{;}                           \CommentTok{// ERROR: PUBLISH\_NO\_PUBLISH\_IDENTITY}
\end{Highlighting}
\end{Shaded}

Omitting a signer is an \textbf{error}, never silent anonymity.
Anonymity must be asked for by name (\texttt{ANONYMOUS}). And the node
key never signs a publish -- even if you wanted it to, you would have to
pass it explicitly as the signer, which is discouraged.

An anonymous message has no provable author, so it carries no
\texttt{signerPubkey}, its Message ID is \texttt{SHA-256(message)}
alone, and it \textbf{cannot be killed} (there is no author to authorize
the retraction).

\subsubsection{5.5 Verifying what you
receive}\label{verifying-what-you-receive}

The network already drops spoofed-signature traffic at a topic's
ingress, but verifying what you consume is good practice if your app
gates on identity:

\begin{Shaded}
\begin{Highlighting}[]
\ImportTok{import}\NormalTok{ \{ verifyEnvelope \} }\ImportTok{from} \StringTok{\textquotesingle{}@axona/protocol\textquotesingle{}}\OperatorTok{;}

\ControlFlowTok{await}\NormalTok{ peer}\OperatorTok{.}\FunctionTok{sub}\NormalTok{(\{ }\DataTypeTok{name}\OperatorTok{:} \StringTok{\textquotesingle{}lobby\textquotesingle{}}\NormalTok{ \}}\OperatorTok{,} \KeywordTok{async}\NormalTok{ (env) }\KeywordTok{=\textgreater{}}\NormalTok{ \{}
  \ControlFlowTok{if}\NormalTok{ (env}\OperatorTok{.}\AttributeTok{signature}\NormalTok{) \{}
    \KeywordTok{const}\NormalTok{ res }\OperatorTok{=} \ControlFlowTok{await} \FunctionTok{verifyEnvelope}\NormalTok{(env)}\OperatorTok{;}   \CommentTok{// returns \{ ok, reason, signed \}}
    \ControlFlowTok{if}\NormalTok{ (}\OperatorTok{!}\NormalTok{res}\OperatorTok{.}\AttributeTok{ok}\NormalTok{) \{ }\BuiltInTok{console}\OperatorTok{.}\FunctionTok{warn}\NormalTok{(}\StringTok{\textquotesingle{}rejected\textquotesingle{}}\OperatorTok{,}\NormalTok{ env}\OperatorTok{.}\AttributeTok{msgId}\OperatorTok{,}\NormalTok{ res}\OperatorTok{.}\AttributeTok{reason}\NormalTok{)}\OperatorTok{;} \ControlFlowTok{return}\OperatorTok{;}\NormalTok{ \}}
\NormalTok{  \}}
  \FunctionTok{appendToFeed}\NormalTok{(env)}\OperatorTok{;}
\NormalTok{\}}\OperatorTok{,}\NormalTok{ \{ }\DataTypeTok{since}\OperatorTok{:} \StringTok{\textquotesingle{}all\textquotesingle{}}\NormalTok{ \})}\OperatorTok{;}
\end{Highlighting}
\end{Shaded}

\begin{quote}
\texttt{verifyEnvelope} returns a \textbf{result object}, not a boolean
-- test \texttt{.ok}. \texttt{if\ (!(await\ verifyEnvelope(env)))} is
always false (an object is truthy), so that guard would silently never
reject anything.
\end{quote}

\texttt{verifyEnvelope} proves the message came from
\texttt{env.signerPubkey} and was not tampered with. It does
\textbf{not} decide whether you \emph{trust} that Author ID -- that is
your application's call.

\begin{center}\rule{0.5\linewidth}{0.5pt}\end{center}

\subsection{6. Topics and addressing}\label{topics-and-addressing}

This section is the home for everything about topic IDs.

\subsubsection{6.1 A topic is a descriptor, not a
string}\label{a-topic-is-a-descriptor-not-a-string}

A topic is a small \textbf{descriptor}, and the Topic ID is a hash of it
with a one-byte region prefix:

\begin{verbatim}
topicId = regionByte || SHA-256(canonical({ owner, name, write }))   // 33 bytes = 66 hex chars
\end{verbatim}

{\def\LTcaptype{none} % do not increment counter
\begin{longtable}[]{@{}
  >{\raggedright\arraybackslash}p{(\linewidth - 2\tabcolsep) * \real{0.5000}}
  >{\raggedright\arraybackslash}p{(\linewidth - 2\tabcolsep) * \real{0.5000}}@{}}
\toprule\noalign{}
\begin{minipage}[b]{\linewidth}\raggedright
Field
\end{minipage} & \begin{minipage}[b]{\linewidth}\raggedright
Meaning
\end{minipage} \\
\midrule\noalign{}
\endhead
\bottomrule\noalign{}
\endlastfoot
\texttt{region} & a real geographic cell; becomes the leading byte of
the ID \\
\texttt{owner} & an \textbf{Author ID} (a publish key), or absent \\
\texttt{name} & the human label -- \texttt{"lobby"},
\texttt{"profile"} \\
\texttt{write} & \texttt{"open"} or \texttt{"owner"} -- \textbf{defaults
by whether \texttt{owner} is set} (6.3) \\
\end{longtable}
}

Because the ID is a pure function of those fields, \textbf{anyone who
knows the fields computes the identical ID} -- no registry, no
coordination.

\subsubsection{\texorpdfstring{6.2 The \texttt{region} field: name or
code, or omit
it}{6.2 The region field: name or code, or omit it}}\label{the-region-field-name-or-code-or-omit-it}

\texttt{region} accepts a region \textbf{name}, \textbf{hex string},
\textbf{decimal string}, or \textbf{number} -- all normalize to the same
region byte:

\begin{Shaded}
\begin{Highlighting}[]
\NormalTok{\{ }\DataTypeTok{region}\OperatorTok{:} \StringTok{\textquotesingle{}useast\textquotesingle{}}\OperatorTok{,} \DataTypeTok{name}\OperatorTok{:} \StringTok{\textquotesingle{}lobby\textquotesingle{}}\NormalTok{ \}   }\CommentTok{// name}
\NormalTok{\{ }\DataTypeTok{region}\OperatorTok{:} \StringTok{\textquotesingle{}0x89\textquotesingle{}}\OperatorTok{,}   \DataTypeTok{name}\OperatorTok{:} \StringTok{\textquotesingle{}lobby\textquotesingle{}}\NormalTok{ \}   }\CommentTok{// hex string   {-}{-}\textbackslash{}  same Topic ID}
\NormalTok{\{ }\DataTypeTok{region}\OperatorTok{:} \DecValTok{137}\OperatorTok{,}      \DataTypeTok{name}\OperatorTok{:} \StringTok{\textquotesingle{}lobby\textquotesingle{}}\NormalTok{ \}   }\CommentTok{// number       {-}{-}/  (leading byte 0x89)}
\end{Highlighting}
\end{Shaded}

\texttt{\textquotesingle{}useast\textquotesingle{}\ ==\ \textquotesingle{}0x89\textquotesingle{}\ ==\ 137}.
\textbf{Omit \texttt{region}} and it defaults to the \emph{publisher's
own node region} (the top byte of its Node ID) -- two co-located peers
converge on the same topic with no region named. To share a topic across
regions, name the region explicitly. An unknown region throws.

Region is \textbf{never} derived from the author key (an Author ID has
no location), and there is \textbf{no global region} -- a topic the
whole world reads picks one real region (a deliberate, app-visible
placement) or is sharded across regions at the app layer.

\subsubsection{\texorpdfstring{6.3 \texttt{write} defaults by whether
there is an
\texttt{owner}}{6.3 write defaults by whether there is an owner}}\label{write-defaults-by-whether-there-is-an-owner}

{\def\LTcaptype{none} % do not increment counter
\begin{longtable}[]{@{}
  >{\raggedright\arraybackslash}p{(\linewidth - 4\tabcolsep) * \real{0.3333}}
  >{\raggedright\arraybackslash}p{(\linewidth - 4\tabcolsep) * \real{0.3333}}
  >{\raggedright\arraybackslash}p{(\linewidth - 4\tabcolsep) * \real{0.3333}}@{}}
\toprule\noalign{}
\begin{minipage}[b]{\linewidth}\raggedright
Descriptor
\end{minipage} & \begin{minipage}[b]{\linewidth}\raggedright
Resolved \texttt{write}
\end{minipage} & \begin{minipage}[b]{\linewidth}\raggedright
Meaning
\end{minipage} \\
\midrule\noalign{}
\endhead
\bottomrule\noalign{}
\endlastfoot
\texttt{\{\ region,\ name\ \}} (no owner) & \texttt{open} & open lobby
-- anyone publishes \\
\texttt{\{\ region,\ owner,\ name\ \}} (no write) &
\textbf{\texttt{owner}} & owned -- only the owner publishes \\
\texttt{\{\ region,\ owner,\ name,\ write:\ \textquotesingle{}owner\textquotesingle{}\ \}}
& \texttt{owner} & same ID as the row above \\
\texttt{\{\ region,\ owner,\ name,\ write:\ \textquotesingle{}open\textquotesingle{}\ \}}
& \texttt{open} & owner-namespaced inbox -- anyone publishes \\
\texttt{\{\ region,\ name,\ write:\ \textquotesingle{}owner\textquotesingle{}\ \}}
(no owner) & \texttt{open} & no owner -\textgreater{} \texttt{write}
ignored \\
\end{longtable}
}

Two rules to remember:

\begin{itemize}
\tightlist
\item
  \textbf{No \texttt{owner} -\textgreater{} the topic is open;
  \texttt{write} is ignored.} You cannot have an owner-only topic with
  no owner.
\item
  \textbf{An \texttt{owner} -\textgreater{} \texttt{write} defaults to
  \texttt{\textquotesingle{}owner\textquotesingle{}}} (the safe
  default). So \texttt{\{owner,\ name\}} and
  \texttt{\{owner,\ name,\ write:\textquotesingle{}owner\textquotesingle{}\}}
  are the \emph{same topic}, and forgetting \texttt{write} can never
  silently make an owned feed world-writable. Pass
  \texttt{write:\textquotesingle{}open\textquotesingle{}} explicitly
  only when you want an owner-namespaced topic anyone may post \emph{to}
  (an inbox).
\end{itemize}

\subsubsection{6.4 The three shapes}\label{the-three-shapes}

\begin{Shaded}
\begin{Highlighting}[]
\ImportTok{import}\NormalTok{ \{ createAuthorIdentity \} }\ImportTok{from} \StringTok{\textquotesingle{}@axona/protocol\textquotesingle{}}\OperatorTok{;}
\KeywordTok{const}\NormalTok{ me }\OperatorTok{=} \ControlFlowTok{await} \FunctionTok{createAuthorIdentity}\NormalTok{(\{ }\DataTypeTok{persistAs}\OperatorTok{:} \StringTok{\textquotesingle{}me\textquotesingle{}}\NormalTok{ \})}\OperatorTok{;}
\end{Highlighting}
\end{Shaded}

\textbf{Open lobby} -- \texttt{\{\ region,\ name\ \}}. Anyone publishes
(each message self-signed or anonymous), anyone subscribes.

\begin{Shaded}
\begin{Highlighting}[]
\ControlFlowTok{await}\NormalTok{ peer}\OperatorTok{.}\FunctionTok{pub}\NormalTok{(\{ }\DataTypeTok{region}\OperatorTok{:} \StringTok{\textquotesingle{}useast\textquotesingle{}}\OperatorTok{,} \DataTypeTok{name}\OperatorTok{:} \StringTok{\textquotesingle{}lobby\textquotesingle{}}\NormalTok{ \}}\OperatorTok{,} \StringTok{\textquotesingle{}hello\textquotesingle{}}\OperatorTok{,}\NormalTok{ \{ }\DataTypeTok{signWith}\OperatorTok{:}\NormalTok{ me \})}\OperatorTok{;}
\ControlFlowTok{await}\NormalTok{ peer}\OperatorTok{.}\FunctionTok{sub}\NormalTok{(\{ }\DataTypeTok{region}\OperatorTok{:} \StringTok{\textquotesingle{}useast\textquotesingle{}}\OperatorTok{,} \DataTypeTok{name}\OperatorTok{:} \StringTok{\textquotesingle{}lobby\textquotesingle{}}\NormalTok{ \}}\OperatorTok{,}\NormalTok{ (env) }\KeywordTok{=\textgreater{}}\NormalTok{ \{}
  \BuiltInTok{console}\OperatorTok{.}\FunctionTok{log}\NormalTok{(env}\OperatorTok{.}\AttributeTok{signerPubkey}\OperatorTok{,}\NormalTok{ env}\OperatorTok{.}\AttributeTok{message}\NormalTok{)}\OperatorTok{;}
\NormalTok{\}}\OperatorTok{,}\NormalTok{ \{ }\DataTypeTok{since}\OperatorTok{:} \StringTok{\textquotesingle{}all\textquotesingle{}}\NormalTok{ \})}\OperatorTok{;}
\end{Highlighting}
\end{Shaded}

\textbf{Author feed (profile / broadcast)} --
\texttt{\{\ region,\ owner,\ name\ \}}, write defaults to
\texttt{owner}. Only messages signed by the owner key are accepted; the
storing node recomputes the ID from the signed descriptor and rejects
anything where \texttt{signer\ !==\ owner}. A profile cannot be spoofed
even by talking directly to a storage node.

\begin{Shaded}
\begin{Highlighting}[]
\KeywordTok{const}\NormalTok{ feed }\OperatorTok{=}\NormalTok{ \{ }\DataTypeTok{region}\OperatorTok{:} \StringTok{\textquotesingle{}useast\textquotesingle{}}\OperatorTok{,} \DataTypeTok{owner}\OperatorTok{:}\NormalTok{ me}\OperatorTok{.}\AttributeTok{authorId}\OperatorTok{,} \DataTypeTok{name}\OperatorTok{:} \StringTok{\textquotesingle{}profile\textquotesingle{}}\NormalTok{ \}}\OperatorTok{;}
\ControlFlowTok{await}\NormalTok{ peer}\OperatorTok{.}\FunctionTok{pub}\NormalTok{(feed}\OperatorTok{,}\NormalTok{ \{ }\DataTypeTok{status}\OperatorTok{:} \StringTok{\textquotesingle{}online\textquotesingle{}}\NormalTok{ \}}\OperatorTok{,}\NormalTok{ \{ }\DataTypeTok{signWith}\OperatorTok{:}\NormalTok{ me \})}\OperatorTok{;}   \CommentTok{// only \textasciigrave{}me\textasciigrave{} may write}
\end{Highlighting}
\end{Shaded}

\textbf{Author inbox (post \emph{to} someone)} --
\texttt{\{\ region,\ owner,\ name,\ write:\ \textquotesingle{}open\textquotesingle{}\ \}}.
The owner namespaces it, but anyone may post (signed as themselves); the
owner reads. ``Reach author X'' = put X's Author ID as \texttt{owner},
sign with your own key.

\begin{Shaded}
\begin{Highlighting}[]
\ControlFlowTok{await}\NormalTok{ peer}\OperatorTok{.}\FunctionTok{pub}\NormalTok{(}
\NormalTok{  \{ }\DataTypeTok{region}\OperatorTok{:} \StringTok{\textquotesingle{}useast\textquotesingle{}}\OperatorTok{,} \DataTypeTok{owner}\OperatorTok{:}\NormalTok{ aliceAuthorId}\OperatorTok{,} \DataTypeTok{name}\OperatorTok{:} \StringTok{\textquotesingle{}inbox\textquotesingle{}}\OperatorTok{,} \DataTypeTok{write}\OperatorTok{:} \StringTok{\textquotesingle{}open\textquotesingle{}}\NormalTok{ \}}\OperatorTok{,}
\NormalTok{  \{ }\DataTypeTok{from}\OperatorTok{:}\NormalTok{ me}\OperatorTok{.}\AttributeTok{authorId}\OperatorTok{,} \DataTypeTok{body}\OperatorTok{:} \StringTok{\textquotesingle{}hi alice\textquotesingle{}}\NormalTok{ \}}\OperatorTok{,}
\NormalTok{  \{ }\DataTypeTok{signWith}\OperatorTok{:}\NormalTok{ me \})}\OperatorTok{;}
\end{Highlighting}
\end{Shaded}

Reading is always open. But the Topic ID is \emph{deterministic, not
secret}: a stranger cannot compute an owned feed's ID without the
owner's Author ID, so the owner shares the ID (or the descriptor) with
readers.

\subsubsection{6.5 Generating a Topic ID to
share}\label{generating-a-topic-id-to-share}

\texttt{deriveTopicId(descriptor)} returns the 66-hex Topic ID -- the
same function the peer uses internally. It works for every topic kind.
This is the \textbf{read / subscribe handle} you hand to someone:

\begin{Shaded}
\begin{Highlighting}[]
\ImportTok{import}\NormalTok{ \{ deriveTopicId \} }\ImportTok{from} \StringTok{\textquotesingle{}@axona/protocol\textquotesingle{}}\OperatorTok{;}

\KeywordTok{const}\NormalTok{ lobbyId }\OperatorTok{=} \ControlFlowTok{await} \FunctionTok{deriveTopicId}\NormalTok{(\{ }\DataTypeTok{region}\OperatorTok{:} \StringTok{\textquotesingle{}useast\textquotesingle{}}\OperatorTok{,} \DataTypeTok{name}\OperatorTok{:} \StringTok{\textquotesingle{}lobby\textquotesingle{}}\NormalTok{ \})}\OperatorTok{;}
\KeywordTok{const}\NormalTok{ feedId  }\OperatorTok{=} \ControlFlowTok{await} \FunctionTok{deriveTopicId}\NormalTok{(\{ }\DataTypeTok{region}\OperatorTok{:} \StringTok{\textquotesingle{}useast\textquotesingle{}}\OperatorTok{,} \DataTypeTok{owner}\OperatorTok{:}\NormalTok{ me}\OperatorTok{.}\AttributeTok{authorId}\OperatorTok{,} \DataTypeTok{name}\OperatorTok{:} \StringTok{\textquotesingle{}profile\textquotesingle{}}\NormalTok{ \})}\OperatorTok{;}
\CommentTok{// {-}\textgreater{} e.g. "89a1b2c3..."  (66 hex; leading "89" is the region byte)}
\CommentTok{// share lobbyId / feedId out{-}of{-}band (URL, QR, message)}
\end{Highlighting}
\end{Shaded}

\texttt{sub}, \texttt{pull}, and \texttt{metrics} accept \textbf{either}
a descriptor object \textbf{or} a bare 66-hex Topic ID:

\begin{Shaded}
\begin{Highlighting}[]
\ControlFlowTok{await}\NormalTok{ peer}\OperatorTok{.}\FunctionTok{sub}\NormalTok{(lobbyId}\OperatorTok{,}\NormalTok{ onMsg}\OperatorTok{,}\NormalTok{ \{ }\DataTypeTok{since}\OperatorTok{:} \StringTok{\textquotesingle{}all\textquotesingle{}}\NormalTok{ \})}\OperatorTok{;}   \CommentTok{// subscribe by ID}
\ControlFlowTok{await}\NormalTok{ peer}\OperatorTok{.}\FunctionTok{pull}\NormalTok{(someMsgId}\OperatorTok{,}\NormalTok{ \{ }\DataTypeTok{topic}\OperatorTok{:}\NormalTok{ feedId \})}\OperatorTok{;}      \CommentTok{// pull by ID}
\ControlFlowTok{await}\NormalTok{ peer}\OperatorTok{.}\FunctionTok{metrics}\NormalTok{(feedId)}\OperatorTok{;}                          \CommentTok{// metrics by ID}
\end{Highlighting}
\end{Shaded}

\subsubsection{6.6 Publishing needs the descriptor, not the
ID}\label{publishing-needs-the-descriptor-not-the-id}

\texttt{pub} (and the owner op \texttt{kill}) requires the
\textbf{descriptor} -- passing a bare ID throws
\texttt{PUBLISH\_INVALID\_TOPIC}. This is by design:

\begin{itemize}
\tightlist
\item
  The storing node must recompute the ID from
  \texttt{\{\ region,\ owner,\ name,\ write\ \}} to enforce the write
  policy (\texttt{signer\ ===\ owner}). A Topic ID is a hash; it cannot
  reveal its \texttt{owner}, so the ID alone cannot prove you are
  authorized to write. (If it could, anyone who learned an owned feed's
  ID could post to it.)
\item
  For an open topic the publish descriptor is just
  \texttt{\{\ region,\ name\ \}} -- tiny and human-meaningful, so
  ``share the topic so someone can post'' means sharing those two
  fields. The owner of an owned topic already holds the full descriptor.
\end{itemize}

\textbf{In short: share the ID for reading; share the descriptor for
writing.}

\subsubsection{6.7 Multiple publishers on one owned topic -- proposed,
not
built}\label{multiple-publishers-on-one-owned-topic-proposed-not-built}

Authorizing additional publishers on an owned topic is \textbf{not
implemented}: \texttt{write:\ \textquotesingle{}owner\textquotesingle{}}
means exactly the one \texttt{owner} key. The roadmap design is an
owner-signed writer roster (the topic still commits to one
\texttt{owner} in its ID; the owner publishes a signed
\texttt{\{\ topicId,\ writers,\ seq\ \}} record; storing nodes accept a
publish if the signer is the owner or is in the latest roster, and the
Topic ID never changes). Do not build against it yet.

\subsubsection{6.8 Quick reference}\label{quick-reference}

{\def\LTcaptype{none} % do not increment counter
\begin{longtable}[]{@{}
  >{\raggedright\arraybackslash}p{(\linewidth - 4\tabcolsep) * \real{0.3333}}
  >{\raggedright\arraybackslash}p{(\linewidth - 4\tabcolsep) * \real{0.3333}}
  >{\raggedright\arraybackslash}p{(\linewidth - 4\tabcolsep) * \real{0.3333}}@{}}
\toprule\noalign{}
\begin{minipage}[b]{\linewidth}\raggedright
\end{minipage} & \begin{minipage}[b]{\linewidth}\raggedright
Open topic
\end{minipage} & \begin{minipage}[b]{\linewidth}\raggedright
Owned topic
\end{minipage} \\
\midrule\noalign{}
\endhead
\bottomrule\noalign{}
\endlastfoot
Descriptor & \texttt{\{\ region,\ name\ \}} &
\texttt{\{\ region,\ owner,\ name\ \}} (write -\textgreater{}
\texttt{owner}) \\
Who may publish & anyone (self-signed / anonymous) & only the
\texttt{owner} key (node-enforced) \\
Who may subscribe & anyone & anyone given the ID / descriptor \\
\texttt{region} forms & name,
\texttt{\textquotesingle{}0x89\textquotesingle{}}, \texttt{137}, or omit
-\textgreater{} node region & same \\
Generate ID to share & \texttt{deriveTopicId(\{\ region,\ name\ \})} &
\texttt{deriveTopicId(\{\ region,\ owner,\ name\ \})} \\
Read by ID & \texttt{sub}/\texttt{pull}/\texttt{metrics} accept the
66-hex ID & same \\
Publish by ID & no -- needs \texttt{\{\ region,\ name\ \}} & no -- needs
the descriptor (write policy) \\
Multiple writers & n/a & proposed (owner-signed roster), not built \\
\end{longtable}
}

\begin{center}\rule{0.5\linewidth}{0.5pt}\end{center}

\subsection{7. Publish and subscribe}\label{publish-and-subscribe}

\subsubsection{7.1 Publish}\label{publish}

\begin{Shaded}
\begin{Highlighting}[]
\KeywordTok{const}\NormalTok{ msgId }\OperatorTok{=} \ControlFlowTok{await}\NormalTok{ peer}\OperatorTok{.}\FunctionTok{pub}\NormalTok{(topicDescriptor}\OperatorTok{,}\NormalTok{ message}\OperatorTok{,}\NormalTok{ \{ signWith \})}\OperatorTok{;}
\CommentTok{// returns the envelope\textquotesingle{}s msgId (64{-}char hex content hash)}
\end{Highlighting}
\end{Shaded}

\begin{itemize}
\tightlist
\item
  \texttt{topicDescriptor} is
  \texttt{\{\ region?,\ owner?,\ name,\ write?\ \}} -- a bare ID is
  rejected (6.6).
\item
  \texttt{message} is any JSON-serializable value.
\item
  \texttt{signWith} is an author identity, or the \texttt{ANONYMOUS}
  sentinel. \textbf{Required} (5.4).
\end{itemize}

A publish above the WebRTC-interoperable floor (\textbf{16 KiB} on the
serialized envelope) throws \texttt{PUBLISH\_PAYLOAD\_TOO\_LARGE}
immediately rather than being silently dropped mid-mesh. Chunk large
payloads with \texttt{@axona/protocol/std/chunk}
(\texttt{publishChunkedBytes}), or -- better for images and documents --
publish a small content reference (hash + size + mime) and transfer the
bytes out-of-band. A non-serializable message (circular ref, BigInt)
throws \texttt{PUBLISH\_INVALID\_MESSAGE}.

\begin{quote}
\textbf{\texttt{std/chunk} is reliable by default (kernel ≥ 3.6.0).}
\texttt{publishChunkedBytes} paces the chunks and then \textbf{verifies
what the mesh cached and re-publishes any gaps} -- \texttt{peer.pub} is
fire-and-forget into the transport buffer, so a fast burst would
otherwise drop chunks before they durably cache and a \emph{reload}
subscriber would never reassemble. You no longer pick a throttle value.
The receiver (\texttt{receiveChunkedBytes}) reassembles, or rejects with
the missing indices on timeout -- it never hangs.
\end{quote}

\begin{quote}
\textbf{Delivery, and why you don't ack (4.x).} A publish routes to the
topic's K-closest \textbf{cohort} and is replicated across it, so it
survives the closest node churning out. There is deliberately \textbf{no
delivery ack to the publisher} --- the transport identity is unlinked
from the author identity, so the network can't (and won't) tell you
``who received it where.'' Two automatic mechanisms cover the gaps so
you don't hand-roll retries: a freshly-joined (``cold'') node re-sends
its first publishes a few times over \textasciitilde1 s (the
\textbf{cold-publish burst}, v4.11.0) so a not-yet-warm routing table
doesn't strand them, and a background retry re-sends a recent publish
until the publisher sees its own \texttt{msgId} land as a root. This is
why §13.3's ``publish before the mesh forms'' pitfall is now a soft
edge, not a cliff.
\end{quote}

\subsubsection{7.2 Subscribe}\label{subscribe}

\begin{Shaded}
\begin{Highlighting}[]
\KeywordTok{const}\NormalTok{ sub }\OperatorTok{=} \ControlFlowTok{await}\NormalTok{ peer}\OperatorTok{.}\FunctionTok{sub}\NormalTok{(topic}\OperatorTok{,}\NormalTok{ (envelope) }\KeywordTok{=\textgreater{}}\NormalTok{ \{}
  \ControlFlowTok{if}\NormalTok{ (envelope}\OperatorTok{.}\AttributeTok{deleted}\NormalTok{) \{ ui}\OperatorTok{.}\FunctionTok{remove}\NormalTok{(envelope}\OperatorTok{.}\AttributeTok{msgId}\NormalTok{)}\OperatorTok{;} \ControlFlowTok{return}\OperatorTok{;}\NormalTok{ \}  }\CommentTok{// a kill arrived}
\NormalTok{  ui}\OperatorTok{.}\FunctionTok{render}\NormalTok{(envelope}\OperatorTok{.}\AttributeTok{message}\NormalTok{)}\OperatorTok{;}
\NormalTok{\}}\OperatorTok{,}\NormalTok{ \{ }\DataTypeTok{since}\OperatorTok{:} \StringTok{\textquotesingle{}all\textquotesingle{}}\NormalTok{ \})}\OperatorTok{;}

\CommentTok{// later}
\ControlFlowTok{await}\NormalTok{ sub}\OperatorTok{.}\FunctionTok{stop}\NormalTok{()}\OperatorTok{;}
\end{Highlighting}
\end{Shaded}

\begin{itemize}
\tightlist
\item
  \texttt{topic} is a descriptor \textbf{or} a bare Topic ID (6.5).
\item
  The handler receives the full envelope (3.4), or a delete marker
  \texttt{\{\ msgId,\ topic,\ deleted:\ true\ \}} when a message is
  retracted -- branch on \texttt{envelope.deleted}.
\item
  \texttt{peer.sub} returns a \texttt{Subscription}; call
  \texttt{.stop()} to leave (or \texttt{peer.unsub(topic)} -- 7.6).
\end{itemize}

\subsubsection{\texorpdfstring{7.3 Replay vs live tail
(\texttt{since})}{7.3 Replay vs live tail (since)}}\label{replay-vs-live-tail-since}

{\def\LTcaptype{none} % do not increment counter
\begin{longtable}[]{@{}
  >{\raggedright\arraybackslash}p{(\linewidth - 2\tabcolsep) * \real{0.5000}}
  >{\raggedright\arraybackslash}p{(\linewidth - 2\tabcolsep) * \real{0.5000}}@{}}
\toprule\noalign{}
\begin{minipage}[b]{\linewidth}\raggedright
\texttt{since}
\end{minipage} & \begin{minipage}[b]{\linewidth}\raggedright
Behavior
\end{minipage} \\
\midrule\noalign{}
\endhead
\bottomrule\noalign{}
\endlastfoot
omitted & Live tail only -- no cached messages. \\
\texttt{\textquotesingle{}all\textquotesingle{}} & Replay everything in
the axons' caches, then live tail. \\
\texttt{\textquotesingle{}latest\textquotesingle{}} & Replay the single
newest retained message (the current value, regardless of age), then
live tail. \\
\texttt{\textless{}number\textgreater{}} & Replay messages newer than
this ms-epoch timestamp, then live tail. \\
\end{longtable}
}

For chat / feed UX you almost always want
\texttt{since:\ \textquotesingle{}all\textquotesingle{}} -- new
subscribers expect history.

\subsubsection{7.4 Queue size and hold
time}\label{queue-size-and-hold-time}

Each topic's replay queue is bounded:

\begin{itemize}
\tightlist
\item
  \textbf{Max messages} -- default \textbf{1024}, with a \textbf{16 MB}
  per-topic byte cap; eviction honors whichever binds. Order comes from
  the signed \texttt{seq}, so every replica evicts the same message and
  caches stay consistent.
\item
  \textbf{Hold time} -- a message past its hold expires and stops being
  delivered or pulled. A \texttt{pull} or \texttt{touch} slides the hold
  forward (bounded by an absolute ceiling), so reads keep a hot message
  alive but cannot pin it forever.
\item
  \textbf{Per-publisher quota (open topics)} -- on an anyone-may-publish
  topic one publisher can hold at most a bounded fraction of the queue,
  so a single flooder cannot evict everyone else. Owner-only topics are
  not quota'd.
\end{itemize}

The cache is in-memory and vanishes when an axon disconnects. A topic
with no subscribers and no host for a short grace window is swept -- so
without a hosting node (7.9), replay history can be lost even before its
hold elapses.

\subsubsection{7.5 Subscriptions are soft-state
leases}\label{subscriptions-are-soft-state-leases}

A subscription is \emph{soft state}. You call \texttt{peer.sub()}
\textbf{once}; the kernel re-announces it to the topic's current axons
on a \textasciitilde10s tick, and an axon forgets a subscriber it has
not heard from in \textasciitilde30s. The 3x headroom means a transient
hiccup does not evict you, and a peer that goes silent and comes back
re-subscribes automatically on the next refresh (with
\texttt{since}-replay backfilling anything still in the hold window).

\textbf{For your application: nothing to manage.} You never renew
subscriptions yourself.

\subsubsection{7.6 Lifecycle: unsub, kill}\label{lifecycle-unsub-kill}

Both are self-authenticating -- authority is proven by the same author
key, no gatekeeper. \texttt{kill} names its author via
\texttt{signWith}.

\begin{Shaded}
\begin{Highlighting}[]
\CommentTok{// Leave a topic (stops all your local subs for it; sends the network unsub).}
\ControlFlowTok{await}\NormalTok{ peer}\OperatorTok{.}\FunctionTok{unsub}\NormalTok{(\{ }\DataTypeTok{region}\OperatorTok{:} \StringTok{\textquotesingle{}useast\textquotesingle{}}\OperatorTok{,} \DataTypeTok{name}\OperatorTok{:} \StringTok{\textquotesingle{}lobby\textquotesingle{}}\NormalTok{ \})}\OperatorTok{;}

\CommentTok{// Retract a message YOU authored ("unsend"). Accepted only if signed by the}
\CommentTok{// SAME author key that signed the original.}
\ControlFlowTok{await}\NormalTok{ peer}\OperatorTok{.}\FunctionTok{kill}\NormalTok{(\{ }\DataTypeTok{region}\OperatorTok{:} \StringTok{\textquotesingle{}useast\textquotesingle{}}\OperatorTok{,} \DataTypeTok{name}\OperatorTok{:} \StringTok{\textquotesingle{}lobby\textquotesingle{}}\NormalTok{ \}}\OperatorTok{,}\NormalTok{ msgId}\OperatorTok{,}\NormalTok{ \{ }\DataTypeTok{signWith}\OperatorTok{:}\NormalTok{ me \})}\OperatorTok{;}
\end{Highlighting}
\end{Shaded}

\begin{quote}
\textbf{Removed in v4.3.0: \texttt{touch} and \texttt{unpub}.}
\texttt{kill} is now the single retraction primitive. \texttt{unpub}
(clearing a whole owned-topic queue) is gone --- a queue empties
naturally as messages age out at the 48 h ceiling. \texttt{touch} (a
hold-time keep-alive) is gone too --- keep a still-relevant message
current by \textbf{re-publishing} it (an upsert that resets the hold
\emph{and} the ceiling; see §7.7). \texttt{peer.touch()} remains
callable as a no-op for source compatibility.
\end{quote}

Honesty note: \texttt{kill} is \textbf{best-effort redaction, not a
cryptographic un-send} -- a subscriber that already has the plaintext
can keep it, and an offline one may never see the purge. An anonymous
message cannot be killed.

\subsubsection{7.7 Message hold time: re-publish refresh, touch, and the
ceiling}\label{message-hold-time-re-publish-refresh-touch-and-the-ceiling}

A message's hold (expiry) is set off its signed \texttt{ts} when a root
stores it, clamped to receive-time and bounded by a 48h ceiling. Three
things to know:

\textbf{Re-publishing the same message refreshes it (an upsert).} The
msgId is \texttt{SHA-256(author\ +\ message)} (no
\texttt{ts}/\texttt{seq}/nonce -- see 3.4), so re-publishing the
identical payload yields the \textbf{same msgId}. As of kernel v3.3.0
the storing node \textbf{upserts}: it replaces the older copy with the
new one, so there's always exactly one entry per msgId -- the newest --
and the new publish sets a \textbf{fresh hold and a fresh 48h ceiling}.
So re-publishing \emph{does} keep a message alive (indefinitely, if you
keep doing it). It is \textbf{delivered once}: subscribers who already
received that msgId are not re-notified, and late subscribers get it
once via replay-on-subscribe. (Different authors of identical text have
\textbf{different} msgIds and are not collapsed; add a nonce to
\texttt{message} for an independently-addressable copy by the same
author.)

\textbf{\texttt{pull} also refreshes, but cannot reset the ceiling.} A
\texttt{pull} slides the hold forward as a side effect of the read
(locally, on the serving replica), clamped to the entry's
\textbf{current} \texttt{ceilingAt}. So a pulled message lives a little
longer without re-notifying; re-publish when you want a fresh ceiling
and to re-establish the content as the newest entry. (Pre-v4.3.0
\texttt{touch} did this deliberately; it's removed --- re-publish
instead.)

\textbf{The ceiling.} An entry can never outlive its
\texttt{ceilingAt\ =\ ts\ +\ 48h} via \texttt{pull}. A re-publish sets a
new \texttt{ts}, hence a new ceiling -- which is why re-publishing can
extend indefinitely while a pull-refresh cannot.

With \texttt{kill}: because there is one entry per msgId, a single
\texttt{kill(topic,\ msgId,\ \{\ signWith\ \})} retracts it -- no stale
copy is left behind. The tombstone that blocks re-acceptance is
finite-lived, and since the msgId is topic-independent, the same content
published to two different topics must be killed in each.

\subsubsection{7.8 Pull and metrics}\label{pull-and-metrics}

\begin{Shaded}
\begin{Highlighting}[]
\CommentTok{// Fetch one envelope by its content hash (e.g. from a reshare link), or the}
\CommentTok{// latest with a null msgId. Returns null on a cache miss / expired message.}
\KeywordTok{const}\NormalTok{ env    }\OperatorTok{=} \ControlFlowTok{await}\NormalTok{ peer}\OperatorTok{.}\FunctionTok{pull}\NormalTok{(msgId}\OperatorTok{,}\NormalTok{ \{ }\DataTypeTok{topic}\OperatorTok{:}\NormalTok{ feedId}\OperatorTok{,} \DataTypeTok{timeoutMs}\OperatorTok{:} \DecValTok{1000}\NormalTok{ \})}\OperatorTok{;}
\KeywordTok{const}\NormalTok{ latest }\OperatorTok{=} \ControlFlowTok{await}\NormalTok{ peer}\OperatorTok{.}\FunctionTok{pull}\NormalTok{(}\KeywordTok{null}\OperatorTok{,}\NormalTok{  \{ }\DataTypeTok{topic}\OperatorTok{:}\NormalTok{ feedId \})}\OperatorTok{;}

\CommentTok{// Coarse, best{-}effort topic state. Works for ANY topic — open or owned.}
\KeywordTok{const}\NormalTok{ m }\OperatorTok{=} \ControlFlowTok{await}\NormalTok{ peer}\OperatorTok{.}\FunctionTok{metrics}\NormalTok{(feedId}\OperatorTok{,}\NormalTok{ \{ }\DataTypeTok{timeoutMs}\OperatorTok{:} \DecValTok{1500}\NormalTok{ \})}\OperatorTok{;}
\CommentTok{// \{ current\_count, seq, subscribers, bytes, publishes, ts, signer, cohortSize, stale \}}
\end{Highlighting}
\end{Shaded}

\texttt{pull} is for ``did I miss this one?'', not durable storage.
\textbf{Nearest-replica reads (v4.11.1--v4.11.2):} a pull is answered by
the first replica it reaches --- a cohort member, a child relay, or a
\texttt{host()} node --- not necessarily the root. \texttt{pull(msgId)}
is \emph{exact} (the content hash pins the message);
\textbf{pull-latest} (\texttt{null} msgId) gives whatever newest that
replica holds --- \emph{recent, eventually-consistent} --- which keeps a
heavily-polled ``current value'' read from hammering the root. Need the
strict newest? Pull a specific \texttt{msgId}.

\texttt{metrics()} (v4.3.0) is a one-shot read of the latest
\textbf{published} metric snapshot --- a topic's roots publish a signed
snapshot to \texttt{metricTopic(T)} every \textasciitilde20 s, and
\texttt{metrics()} subscribes to it briefly and aggregates across the
cohort (\texttt{stale:true} if no publisher roots the topic yet). It
works for \textbf{any} topic, open or owned --- an owned topic's metrics
are public too. \texttt{current\_count} is the live retained count;
\texttt{seq} is the dense message counter (total events ever, kills
included); \texttt{subscribers} is the topic-wide total (summed across
the cohort); \texttt{cohortSize} is how many roots reported. For UX, not
billing.

\begin{quote}
\textbf{For a live count, prefer \texttt{sub(metricTopic(T),\ …)}
directly} (next section): one standing subscription, every update pushed
to you, plus a rolling history. \texttt{metrics()} is the convenience
one-shot.
\end{quote}

\subsubsection{\texorpdfstring{7.8b Watching metrics live
(\texttt{metricTopic})}{7.8b Watching metrics live (metricTopic)}}\label{b-watching-metrics-live-metrictopic}

For a continuously-displayed count, \textbf{subscribe to the topic's
metrics.} A topic's primary root publishes a signed snapshot to a
\emph{derived, open} topic every \textasciitilde20 s --- for
\textbf{both open and owned} data topics; you compute that topic with
\texttt{metricTopic()} and \texttt{sub()} it like any other:

\begin{Shaded}
\begin{Highlighting}[]
\ImportTok{import}\NormalTok{ \{ deriveTopicId}\OperatorTok{,}\NormalTok{ metricTopic \} }\ImportTok{from} \StringTok{\textquotesingle{}@axona/protocol\textquotesingle{}}\OperatorTok{;}

\KeywordTok{const}\NormalTok{ lobbyId }\OperatorTok{=} \ControlFlowTok{await} \FunctionTok{deriveTopicId}\NormalTok{(\{ }\DataTypeTok{region}\OperatorTok{:} \StringTok{\textquotesingle{}useast\textquotesingle{}}\OperatorTok{,} \DataTypeTok{name}\OperatorTok{:} \StringTok{\textquotesingle{}lobby\textquotesingle{}}\NormalTok{ \})}\OperatorTok{;}
\ControlFlowTok{await}\NormalTok{ peer}\OperatorTok{.}\FunctionTok{sub}\NormalTok{(}\FunctionTok{metricTopic}\NormalTok{(lobbyId)}\OperatorTok{,}\NormalTok{ (env) }\KeywordTok{=\textgreater{}}\NormalTok{ \{}
  \KeywordTok{const}\NormalTok{ m }\OperatorTok{=} \BuiltInTok{JSON}\OperatorTok{.}\FunctionTok{parse}\NormalTok{(env}\OperatorTok{.}\AttributeTok{message}\NormalTok{)}\OperatorTok{;}   \CommentTok{// \{ topic, ts, by, signer, current\_count, seq, subscribers, bytes \}}
  \FunctionTok{showRoomCount}\NormalTok{(m}\OperatorTok{.}\AttributeTok{subscribers}\OperatorTok{,}\NormalTok{ m}\OperatorTok{.}\AttributeTok{current\_count}\NormalTok{)}\OperatorTok{;}
\NormalTok{\}}\OperatorTok{,}\NormalTok{ \{ }\DataTypeTok{since}\OperatorTok{:} \StringTok{\textquotesingle{}all\textquotesingle{}}\NormalTok{ \})}\OperatorTok{;}                  \CommentTok{// latest snapshot + a rolling \textasciitilde{}48 h trend}
\end{Highlighting}
\end{Shaded}

You pay one subscription instead of a poll, you get every update pushed
to you, and because each snapshot is an ordinary retained message that
ages out at the 48 h hold ceiling,
\texttt{since:\textquotesingle{}all\textquotesingle{}} hands you a
\textbf{rolling \textasciitilde48 h history} --- enough to plot a trend
line, free.

Three things to keep in mind:

\begin{itemize}
\tightlist
\item
  \textbf{Advisory, not authoritative.} The metric topic is \emph{open}
  (anyone can publish to it), so treat a snapshot as a hint. If you only
  trust your own relays, check \texttt{env.signerPubkey} against their
  keys.
\item
  \textbf{Owned topics included (v4.3.0).} An owned
  (\texttt{write:\textquotesingle{}owner\textquotesingle{}}) topic's
  metrics are published to its (open) metric topic too --- so anyone can
  subscribe to an owned topic's activity metrics without owning it. The
  \emph{messages} stay write- gated; only the activity counts are
  public.
\item
  \textbf{You don't publish these --- relays do.} As an app you only
  ever \emph{read} metric topics. The publish side runs on
  infrastructure (see §7.9 / the relay).
\end{itemize}

\subsubsection{\texorpdfstring{7.9 Hosting a topic without subscribing
(\texttt{host} /
\texttt{unhost})}{7.9 Hosting a topic without subscribing (host / unhost)}}\label{hosting-a-topic-without-subscribing-host-unhost}

A relay or archive node wants to \emph{store and serve} a topic without
consuming it. That is \texttt{host}, decoupled from \texttt{sub} on
purpose -- hosting is ``I will serve this for others,'' subscribing is
``I want to receive this.''

\begin{Shaded}
\begin{Highlighting}[]
\ControlFlowTok{await}\NormalTok{ peer}\OperatorTok{.}\FunctionTok{host}\NormalTok{()}\OperatorTok{;}                                  \CommentTok{// host my keyspace neighborhood}
\ControlFlowTok{await}\NormalTok{ peer}\OperatorTok{.}\FunctionTok{host}\NormalTok{(\{ }\DataTypeTok{region}\OperatorTok{:} \StringTok{\textquotesingle{}useast\textquotesingle{}}\OperatorTok{,} \DataTypeTok{name}\OperatorTok{:} \StringTok{\textquotesingle{}lobby\textquotesingle{}}\NormalTok{ \})}\OperatorTok{;} \CommentTok{// host one specific topic}
\ControlFlowTok{await}\NormalTok{ peer}\OperatorTok{.}\FunctionTok{unhost}\NormalTok{(\{ }\DataTypeTok{region}\OperatorTok{:} \StringTok{\textquotesingle{}useast\textquotesingle{}}\OperatorTok{,} \DataTypeTok{name}\OperatorTok{:} \StringTok{\textquotesingle{}lobby\textquotesingle{}}\NormalTok{ \})}\OperatorTok{;}
\ControlFlowTok{await}\NormalTok{ peer}\OperatorTok{.}\FunctionTok{unhost}\NormalTok{()}\OperatorTok{;}                                 \CommentTok{// stop keyspace hosting}
\end{Highlighting}
\end{Shaded}

\texttt{host()} with no argument makes the node a willing root for
whatever topics land near its Node ID -- this is how
\texttt{axona-relay} provides durable history. \texttt{host(topic)} pins
one topic. Hosting registers no handler and delivers nothing to your
app.

\begin{center}\rule{0.5\linewidth}{0.5pt}\end{center}

\subsection{8. Direct messaging}\label{direct-messaging}

Not all peer-to-peer communication is pub/sub. For 1:1 messages,
\texttt{AxonaPeer} exposes a direct API. There is one handler per peer;
dispatch on payload shape if you want multiple message kinds.

\begin{Shaded}
\begin{Highlighting}[]
\CommentTok{// Request/response {-}{-} Alice asks Bob, awaits his handler\textquotesingle{}s return value.}
\KeywordTok{const}\NormalTok{ reply }\OperatorTok{=} \ControlFlowTok{await}\NormalTok{ alice}\OperatorTok{.}\FunctionTok{send}\NormalTok{(bobNodeIdHex}\OperatorTok{,}\NormalTok{ \{ }\DataTypeTok{kind}\OperatorTok{:} \StringTok{\textquotesingle{}ping\textquotesingle{}}\OperatorTok{,} \DataTypeTok{body}\OperatorTok{:} \StringTok{\textquotesingle{}hello\textquotesingle{}}\NormalTok{ \})}\OperatorTok{;}

\CommentTok{// Fire{-}and{-}forget.}
\ControlFlowTok{await}\NormalTok{ alice}\OperatorTok{.}\FunctionTok{notify}\NormalTok{(bobNodeIdHex}\OperatorTok{,}\NormalTok{ \{ }\DataTypeTok{kind}\OperatorTok{:} \StringTok{\textquotesingle{}typing\textquotesingle{}}\OperatorTok{,} \DataTypeTok{user}\OperatorTok{:} \StringTok{\textquotesingle{}alice\textquotesingle{}}\NormalTok{ \})}\OperatorTok{;}

\CommentTok{// Bob\textquotesingle{}s single handler {-}{-} invoked for every inbound direct message.}
\NormalTok{bob}\OperatorTok{.}\FunctionTok{onMessage}\NormalTok{(}\KeywordTok{async}\NormalTok{ (senderId}\OperatorTok{,}\NormalTok{ message) }\KeywordTok{=\textgreater{}}\NormalTok{ \{}
  \ControlFlowTok{switch}\NormalTok{ (message}\OperatorTok{?.}\AttributeTok{kind}\NormalTok{) \{}
    \ControlFlowTok{case} \StringTok{\textquotesingle{}ping\textquotesingle{}}\OperatorTok{:}   \ControlFlowTok{return}\NormalTok{ \{ }\DataTypeTok{reply}\OperatorTok{:} \StringTok{\textquotesingle{}pong\textquotesingle{}}\NormalTok{ \}}\OperatorTok{;}      \CommentTok{// becomes Alice\textquotesingle{}s send() resolution}
    \ControlFlowTok{case} \StringTok{\textquotesingle{}typing\textquotesingle{}}\OperatorTok{:} \BuiltInTok{console}\OperatorTok{.}\FunctionTok{log}\NormalTok{(}\VerbatimStringTok{\textasciigrave{}}\SpecialCharTok{$\{}\NormalTok{senderId}\SpecialCharTok{\}}\VerbatimStringTok{ typing\textasciigrave{}}\NormalTok{)}\OperatorTok{;} \ControlFlowTok{return}\OperatorTok{;}
    \ControlFlowTok{default}\OperatorTok{:}       \ControlFlowTok{return}\NormalTok{ \{ }\DataTypeTok{error}\OperatorTok{:} \StringTok{\textquotesingle{}unknown kind\textquotesingle{}}\NormalTok{ \}}\OperatorTok{;}
\NormalTok{  \}}
\NormalTok{\})}\OperatorTok{;}
\end{Highlighting}
\end{Shaded}

\begin{itemize}
\tightlist
\item
  The target is a \textbf{66-char hex Node ID} (the value
  \texttt{node.id} / \texttt{peer.getNodeId()} returns), and the peer
  must already have a channel to it. Routing to a peer you have no
  channel with is the job of higher layers.
\item
  \texttt{senderId} is set by the transport at the receiving end; trust
  it for routing, but verify application-layer identity yourself if it
  matters.
\item
  \texttt{onMessage} is \textbf{single-handler} -- calling it again
  replaces the previous handler. Dispatch on \texttt{message.kind} for
  multiple kinds.
\item
  \texttt{send}/\texttt{notify} do \textbf{not} sign the way
  \texttt{pub} does. For end-to-end authentic DMs, sign your payload
  yourself and verify on receipt.
\end{itemize}

\begin{center}\rule{0.5\linewidth}{0.5pt}\end{center}

\subsection{9. Mesh introspection}\label{mesh-introspection}

\subsubsection{9.1 Who's online}\label{whos-online}

\begin{Shaded}
\begin{Highlighting}[]
\NormalTok{peer}\OperatorTok{.}\FunctionTok{peers}\NormalTok{()}\OperatorTok{;}   \CommentTok{// {-}\textgreater{} [\textquotesingle{}\textless{}66 hex\textgreater{}\textquotesingle{}, ...] node IDs currently in the synaptome}

\KeywordTok{const}\NormalTok{ off1 }\OperatorTok{=}\NormalTok{ peer}\OperatorTok{.}\FunctionTok{onPeerJoin}\NormalTok{((peerId) }\KeywordTok{=\textgreater{}} \BuiltInTok{console}\OperatorTok{.}\FunctionTok{log}\NormalTok{(}\StringTok{\textquotesingle{}hi\textquotesingle{}}\OperatorTok{,}\NormalTok{ peerId))}\OperatorTok{;}
\KeywordTok{const}\NormalTok{ off2 }\OperatorTok{=}\NormalTok{ peer}\OperatorTok{.}\FunctionTok{onPeerLeave}\NormalTok{((peerId) }\KeywordTok{=\textgreater{}} \BuiltInTok{console}\OperatorTok{.}\FunctionTok{log}\NormalTok{(}\StringTok{\textquotesingle{}bye\textquotesingle{}}\OperatorTok{,}\NormalTok{ peerId))}\OperatorTok{;}
\FunctionTok{off1}\NormalTok{()}\OperatorTok{;} \FunctionTok{off2}\NormalTok{()}\OperatorTok{;}   \CommentTok{// each returns an unsubscribe function}
\end{Highlighting}
\end{Shaded}

\texttt{peers()} is the synaptome -- everyone in a small mesh, a
constant-sized subset in a large one.

\subsubsection{9.2 Lookup}\label{lookup}

\begin{Shaded}
\begin{Highlighting}[]
\KeywordTok{const}\NormalTok{ result }\OperatorTok{=} \ControlFlowTok{await}\NormalTok{ peer}\OperatorTok{.}\FunctionTok{lookup}\NormalTok{(}\BuiltInTok{BigInt}\NormalTok{(}\StringTok{\textquotesingle{}0x\textquotesingle{}} \OperatorTok{+}\NormalTok{ targetHex))}\OperatorTok{;}
\CommentTok{// \{ found, hops, time, path \}}
\end{Highlighting}
\end{Shaded}

An iterative XOR-routed walk; warms the local synaptome as a side
effect.

\subsubsection{9.3 Health}\label{health}

\begin{Shaded}
\begin{Highlighting}[]
\NormalTok{peer}\OperatorTok{.}\FunctionTok{health}\NormalTok{()}\OperatorTok{;}
\CommentTok{// \{}
\CommentTok{//   nodeId, synaptomeSize, peers: [...], subscriptions,}
\CommentTok{//   axonRoles: [\{ topic, isRoot, children, cacheSize \}],}
\CommentTok{//   hosting, wireVersion, started,}
\CommentTok{//   transport: \{ boundCount, meshChannels, meshOpen, meshBound, bridgeState \} | null,}
\CommentTok{//   meshDegraded,   // true =\textgreater{} channels OPEN but NOT yet authenticated}
\CommentTok{// \}}
\end{Highlighting}
\end{Shaded}

\texttt{meshDegraded} is the routing-truth signal: it distinguishes
``channels are open'' from ``channels are open \emph{and} carrying
authenticated routing.'' A single true tick is a normal mid-handshake
transient; a value that stays true across several polls means the mesh
looks connected but is not actually routing. \texttt{boundCount} /
\texttt{meshBound} are the honest usable-peer counts. Cheap; safe to
poll on a UI tick.

\subsubsection{9.4 Logs, errors, upgrades}\label{logs-errors-upgrades}

\begin{Shaded}
\begin{Highlighting}[]
\KeywordTok{const}\NormalTok{ off }\OperatorTok{=}\NormalTok{ peer}\OperatorTok{.}\FunctionTok{onLog}\NormalTok{(}\StringTok{\textquotesingle{}warn\textquotesingle{}}\OperatorTok{,}\NormalTok{ (msg}\OperatorTok{,}\NormalTok{ ctx) }\KeywordTok{=\textgreater{}}\NormalTok{ \{ }\CommentTok{/* \textquotesingle{}debug\textquotesingle{}|\textquotesingle{}info\textquotesingle{}|\textquotesingle{}warn\textquotesingle{}|\textquotesingle{}error\textquotesingle{} */}\NormalTok{ \})}\OperatorTok{;}
\NormalTok{peer}\OperatorTok{.}\FunctionTok{onError}\NormalTok{((err) }\KeywordTok{=\textgreater{}}\NormalTok{ \{ }\CommentTok{/* background AxonaError */}\NormalTok{ \})}\OperatorTok{;}
\NormalTok{peer}\OperatorTok{.}\FunctionTok{onUpgradeRequired}\NormalTok{((err) }\KeywordTok{=\textgreater{}}\NormalTok{ \{ }\CommentTok{/* show an upgrade banner */}\NormalTok{ \})}\OperatorTok{;}
\end{Highlighting}
\end{Shaded}

\texttt{onLog} takes the \textbf{level first}, then the handler, and
returns an unsubscribe. The kernel forwards its security drop-path logs
(bad signature, stale, oversize, write-policy violation, unauthorized
kill, \ldots) to \texttt{onLog} so you can surface them.

\begin{center}\rule{0.5\linewidth}{0.5pt}\end{center}

\subsection{10. Persistence}\label{persistence}

\subsubsection{10.1 The author key (the one you usually
persist)}\label{the-author-key-the-one-you-usually-persist}

For most apps the only persistence decision is whether your author is
durable.
\texttt{createAuthorIdentity(\{\ persistAs:\ \textquotesingle{}me\textquotesingle{}\ \})}
(5.2) handles it -- no adapter, no manual dump/load.

\subsubsection{10.2 The PersistenceAdapter (peer
state)}\label{the-persistenceadapter-peer-state}

To persist \emph{peer} state -- the node identity envelope, the
synaptome, the active subscription list -- pass a
\texttt{PersistenceAdapter} as the constructor's \texttt{persist}
option. The kernel ships two, behind sub-path imports so neither pulls
\texttt{indexedDB} into Node nor \texttt{node:fs} into the browser
bundle:

\begin{Shaded}
\begin{Highlighting}[]
\ImportTok{import}\NormalTok{ \{ IndexedDBPersistence \} }\ImportTok{from} \StringTok{\textquotesingle{}@axona/protocol/persistence/indexeddb.js\textquotesingle{}}\OperatorTok{;}
\CommentTok{// import \{ FilePersistence \} from \textquotesingle{}@axona/protocol/persistence/file.js\textquotesingle{}; // Node}

\KeywordTok{const}\NormalTok{ peer }\OperatorTok{=} \KeywordTok{new} \FunctionTok{AxonaPeer}\NormalTok{(\{}
\NormalTok{  domain}\OperatorTok{,}\NormalTok{ node}\OperatorTok{,}\NormalTok{ nodeIdentity}\OperatorTok{,}\NormalTok{ transport}\OperatorTok{,}
  \DataTypeTok{persist}\OperatorTok{:} \KeywordTok{new} \FunctionTok{IndexedDBPersistence}\NormalTok{(\{ }\DataTypeTok{dbName}\OperatorTok{:} \StringTok{\textquotesingle{}my{-}app\textquotesingle{}}\NormalTok{ \})}\OperatorTok{,}
\NormalTok{\})}\OperatorTok{;}
\ControlFlowTok{await}\NormalTok{ peer}\OperatorTok{.}\FunctionTok{start}\NormalTok{()}\OperatorTok{;}
\end{Highlighting}
\end{Shaded}

On \texttt{start()} the peer loads the node identity (if you did not
pass one), the synaptome seed, and the subscription list (exposed as
\texttt{peer.pendingSubscriptions} for you to re-register handlers --
functions do not serialize). Writes are debounced; \texttt{leave()}
force-flushes.

\subsubsection{10.3 The node identity envelope
(rarely)}\label{the-node-identity-envelope-rarely}

If you want a stable Node ID across restarts (rare -- see 4.2),
round-trip the node identity through \texttt{dumpIdentity} /
\texttt{loadIdentity}:

\begin{Shaded}
\begin{Highlighting}[]
\ImportTok{import}\NormalTok{ \{ dumpIdentity}\OperatorTok{,}\NormalTok{ loadIdentity \} }\ImportTok{from} \StringTok{\textquotesingle{}@axona/protocol\textquotesingle{}}\OperatorTok{;}

\NormalTok{localStorage}\OperatorTok{.}\FunctionTok{setItem}\NormalTok{(}\StringTok{\textquotesingle{}node\textquotesingle{}}\OperatorTok{,} \BuiltInTok{JSON}\OperatorTok{.}\FunctionTok{stringify}\NormalTok{(}\ControlFlowTok{await} \FunctionTok{dumpIdentity}\NormalTok{(nodeIdentity)))}\OperatorTok{;}
\KeywordTok{const}\NormalTok{ restored }\OperatorTok{=} \ControlFlowTok{await} \FunctionTok{loadIdentity}\NormalTok{(}\BuiltInTok{JSON}\OperatorTok{.}\FunctionTok{parse}\NormalTok{(localStorage}\OperatorTok{.}\FunctionTok{getItem}\NormalTok{(}\StringTok{\textquotesingle{}node\textquotesingle{}}\NormalTok{)))}\OperatorTok{;}
\end{Highlighting}
\end{Shaded}

\texttt{loadIdentity} verifies the stored Node ID is internally
consistent and that the private key actually corresponds to the public
key, throwing on corruption.

\subsubsection{10.4 Manual snapshots}\label{manual-snapshots}

\begin{Shaded}
\begin{Highlighting}[]
\KeywordTok{const}\NormalTok{ state    }\OperatorTok{=} \ControlFlowTok{await}\NormalTok{ peer}\OperatorTok{.}\FunctionTok{snapshot}\NormalTok{()}\OperatorTok{;}   \CommentTok{// \{ identity?, subscriptions, synaptome, ... \}}
\KeywordTok{const}\NormalTok{ restored }\OperatorTok{=} \ControlFlowTok{await}\NormalTok{ AxonaPeer}\OperatorTok{.}\FunctionTok{fromSnapshot}\NormalTok{(state}\OperatorTok{,}\NormalTok{ \{ engine}\OperatorTok{,}\NormalTok{ node}\OperatorTok{,}\NormalTok{ transport \})}\OperatorTok{;}
\end{Highlighting}
\end{Shaded}

\texttt{fromSnapshot} is a static factory; the restored peer is
constructed and pre-loaded but not connected -- call
\texttt{peer.join()} to bring the transport up.

\begin{center}\rule{0.5\linewidth}{0.5pt}\end{center}

\subsection{11. The bridge (Docker
production)}\label{the-bridge-docker-production}

The bridge does three things:

\begin{enumerate}
\def\labelenumi{\arabic{enumi}.}
\tightlist
\item
  \textbf{Signaling for WebRTC} -- peers exchange SDP offers/answers/ICE
  through it on a per-pair basis.
\item
  \textbf{Peer-list distribution} -- a connecting peer gets the current
  peer list so it can initiate WebRTC offers.
\item
  \textbf{Universal-hub peer} -- the bridge runs its own peer, meshed to
  every connected browser, acting as a routing hop for peers that cannot
  reach each other directly.
\end{enumerate}

\subsubsection{11.1 Using the public
bridges}\label{using-the-public-bridges}

Point at the network for the line you're building on:

\begin{verbatim}
wss://testnet.axona.net        # the 4.x line (this guide) — kernel 4.11.2
wss://bridge.axona.net         # production east — still the 3.x line
wss://bridge-west.axona.net    # production west — still the 3.x line
\end{verbatim}

The 4.x line runs on \textbf{testnet}; the production federated pair is
still on 3.x, and the wire major partitions them (a 4.x peer and a 3.x
bridge reject each other at the handshake). Open a WebSocket (the
\texttt{webTransport} factory does the handshake for you) and you are on
the network. A bridge advertises itself in the public bridge directory
so clients can discover and fail over between bridges.

\subsubsection{11.2 Running your own bridge (Docker
stack)}\label{running-your-own-bridge-docker-stack}

Production runs a \textbf{three-container Docker stack} -- this is the
current deployment, not the old \texttt{npm\ start} + systemd path:

\begin{itemize}
\tightlist
\item
  \textbf{\texttt{bridge}} -- the Node bridge (internal only; not
  published directly).
\item
  \textbf{\texttt{caddy}} -- a reverse proxy with \textbf{automatic
  HTTPS} (Let's Encrypt) that terminates TLS and proxies
  \texttt{wss://\$DOMAIN} to the bridge.
\item
  \textbf{\texttt{coturn}} -- a TURN relay on host networking (it needs
  the full UDP relay port range and the real client source IPs, which
  Docker's bridge network would hide). Always deploy TURN with a bridge
  -- a STUN-only bridge cannot relay for peers behind symmetric NAT.
\end{itemize}

\begin{verbatim}
git clone https://github.com/axona-net/axona-bridge
cd axona-bridge
cp .env.docker.example .env
$EDITOR .env            # set DOMAIN, PUBLIC_IP, TURN_AUTH_SECRET (and BRIDGE_PUBLIC_URL)
docker compose up -d --build
curl https://$DOMAIN/healthz
\end{verbatim}

Prerequisites: a Linux Docker host with ports 80/443 free and DNS
pointing at it. The full recipe (including the
systemd/\texttt{install.sh} alternative, which is the simpler path if
you need \texttt{turns://} TLS on the TURN port) is in
\texttt{axona-bridge/deploy/README.md}.

\subsubsection{11.3 Bridge environment (the ones that
matter)}\label{bridge-environment-the-ones-that-matter}

{\def\LTcaptype{none} % do not increment counter
\begin{longtable}[]{@{}
  >{\raggedright\arraybackslash}p{(\linewidth - 2\tabcolsep) * \real{0.5000}}
  >{\raggedright\arraybackslash}p{(\linewidth - 2\tabcolsep) * \real{0.5000}}@{}}
\toprule\noalign{}
\begin{minipage}[b]{\linewidth}\raggedright
Var
\end{minipage} & \begin{minipage}[b]{\linewidth}\raggedright
Purpose
\end{minipage} \\
\midrule\noalign{}
\endhead
\bottomrule\noalign{}
\endlastfoot
\texttt{DOMAIN} & The bridge's public hostname (Caddy provisions a cert
for it). \\
\texttt{PUBLIC\_IP} & The host's public IPv4, for coturn's
\texttt{-\/-external-ip}. \\
\texttt{TURN\_AUTH\_SECRET} & Shared static-auth secret; the bridge
mints time-bound TURN creds with it and coturn validates against it.
Must match. \\
\texttt{BRIDGE\_PUBLIC\_URL} &
\texttt{wss://\textless{}this-bridge\textgreater{}} so it advertises
itself in the bridge directory. \\
\texttt{BRIDGE\_DIRECTORY} & \texttt{off} to opt a private/testnet
bridge out of the public directory. \\
\texttt{BRIDGE\_UPSTREAMS} & Comma-separated bridges to federate into
(defaults to the public east/west pair). \\
\texttt{BRIDGE\_LAT} / \texttt{BRIDGE\_LNG} /
\texttt{BRIDGE\_REGION\_LABEL} & The bridge's own region. \\
\texttt{PORT} / \texttt{LOG\_LEVEL} & Internal port (8080) and log
level. \\
\end{longtable}
}

A persistent \texttt{bridge-data} volume holds the bridge's identity
keypair, so its Node ID survives restarts.

\subsubsection{11.4 Federation}\label{federation}

When \texttt{BRIDGE\_DIRECTORY} is on, a bridge also bootstraps
\emph{into} the live mesh as a node (an outbound uplink to a known
bridge) so its directory entry is visible network-wide.
\texttt{bridge.axona.net} and \texttt{bridge-west.axona.net} are
federated this way. \texttt{/healthz} exposes
\texttt{uplink.\{\ upstream,\ connected\ \}} and
\texttt{directory.\{\ enabled,\ url\ \}} so you can confirm what is
live.

\subsubsection{11.5 Surface the kernel
version}\label{surface-the-kernel-version}

Every Axona deployment should expose \texttt{KERNEL\_VERSION} at a
runtime-inspectable surface so it is obvious which kernel is live: the
bridge at \texttt{/healthz}, an app in its version row, a Node peer in a
startup log. Import the constant rather than hard-coding a string.

\begin{center}\rule{0.5\linewidth}{0.5pt}\end{center}

\subsection{12. Worked example: a regional chat
app}\label{worked-example-a-regional-chat-app}

Let's build \textbf{AxonaTalk}: users join named rooms in a shared
region, see history, and can DM each other. This follows the real
pattern in \texttt{apps/axona-minimal/}.

\subsubsection{12.1 The plan}\label{the-plan}

\begin{itemize}
\tightlist
\item
  \textbf{Rooms} are \textbf{open topics} pinned to one region, so every
  participant derives the same Topic ID regardless of where they sit:
  \texttt{\{\ region:\ \textquotesingle{}useast\textquotesingle{},\ name:\ room\ \}}.
\item
  \textbf{Each user} has a \textbf{durable author} (\texttt{persistAs})
  so their authorship is stable across reloads, and an \textbf{ephemeral
  node identity} rooted at their real location.
\item
  The protocol never reveals where a sender is. If we want to
  \emph{show} a sender's region, the app voluntarily includes its own
  Node ID in the payload -- an app choice, not a protocol disclosure.
\item
  \textbf{DMs} ride \texttt{peer.send}, not pub/sub.
\end{itemize}

\subsubsection{12.2 Message shape}\label{message-shape}

\begin{Shaded}
\begin{Highlighting}[]
\CommentTok{// A chat message (the \textasciigrave{}message\textasciigrave{} field of the envelope):}
\NormalTok{\{ }\DataTypeTok{text}\OperatorTok{:} \StringTok{\textquotesingle{}hi\textquotesingle{}}\OperatorTok{,} \DataTypeTok{node}\OperatorTok{:} \StringTok{\textquotesingle{}\textless{}my Node ID hex\textgreater{}\textquotesingle{}}\NormalTok{ \}   }\CommentTok{// \textasciigrave{}node\textasciigrave{} is voluntary region disclosure}

\CommentTok{// A DM (sent via peer.send):}
\NormalTok{\{ }\DataTypeTok{kind}\OperatorTok{:} \StringTok{\textquotesingle{}axonatalk:dm\textquotesingle{}}\OperatorTok{,} \DataTypeTok{from}\OperatorTok{:} \StringTok{\textquotesingle{}\textless{}my Author ID\textgreater{}\textquotesingle{}}\OperatorTok{,} \DataTypeTok{text}\OperatorTok{:} \StringTok{\textquotesingle{}...\textquotesingle{}}\OperatorTok{,} \DataTypeTok{ts}\OperatorTok{:} \FloatTok{1779.}\OperatorTok{..}\NormalTok{ \}}
\end{Highlighting}
\end{Shaded}

\subsubsection{12.3 Connect}\label{connect}

\begin{Shaded}
\begin{Highlighting}[]
\ImportTok{import}\NormalTok{ \{}
\NormalTok{  AxonaPeer}\OperatorTok{,}\NormalTok{ AxonaDomain}\OperatorTok{,}\NormalTok{ NeuronNode}\OperatorTok{,}
\NormalTok{  createNodeIdentity}\OperatorTok{,}\NormalTok{ createAuthorIdentity}\OperatorTok{,}
\NormalTok{  deriveTopicId}\OperatorTok{,}\NormalTok{ regionName}\OperatorTok{,}\NormalTok{ regionCenter}\OperatorTok{,}\NormalTok{ KERNEL\_VERSION}\OperatorTok{,}
\NormalTok{\} }\ImportTok{from} \StringTok{\textquotesingle{}@axona/protocol\textquotesingle{}}\OperatorTok{;}
\ImportTok{import}\NormalTok{ \{ webTransport \} }\ImportTok{from} \StringTok{\textquotesingle{}@axona/protocol/transport/web/index.js\textquotesingle{}}\OperatorTok{;}

\KeywordTok{const}\NormalTok{ BRIDGE       }\OperatorTok{=} \StringTok{\textquotesingle{}wss://testnet.axona.net\textquotesingle{}}\OperatorTok{;}  \CommentTok{// the 4.x line runs on testnet (prod is 3.x)}
\KeywordTok{const}\NormalTok{ TOPIC\_REGION }\OperatorTok{=} \StringTok{\textquotesingle{}useast\textquotesingle{}}\OperatorTok{;}                   \CommentTok{// rooms are pinned here for everyone}

\KeywordTok{let}\NormalTok{ peer}\OperatorTok{,}\NormalTok{ nodeIdentity}\OperatorTok{,}\NormalTok{ author}\OperatorTok{;}

\CommentTok{// Real geolocation {-}\textgreater{} the user\textquotesingle{}s actual S2 cell; denied {-}\textgreater{} the room region.}
\KeywordTok{function} \FunctionTok{whereAmI}\NormalTok{() \{}
  \KeywordTok{const}\NormalTok{ fallback }\OperatorTok{=} \FunctionTok{regionCenter}\NormalTok{(TOPIC\_REGION)}\OperatorTok{;}
  \ControlFlowTok{if}\NormalTok{ (}\OperatorTok{!}\BuiltInTok{navigator}\OperatorTok{.}\AttributeTok{geolocation}\NormalTok{) }\ControlFlowTok{return} \BuiltInTok{Promise}\OperatorTok{.}\FunctionTok{resolve}\NormalTok{(fallback)}\OperatorTok{;}
  \ControlFlowTok{return} \KeywordTok{new} \BuiltInTok{Promise}\NormalTok{((resolve) }\KeywordTok{=\textgreater{}}\NormalTok{ \{}
    \BuiltInTok{navigator}\OperatorTok{.}\AttributeTok{geolocation}\OperatorTok{.}\FunctionTok{getCurrentPosition}\NormalTok{(}
\NormalTok{      (p)  }\KeywordTok{=\textgreater{}} \FunctionTok{resolve}\NormalTok{(\{ }\DataTypeTok{lat}\OperatorTok{:}\NormalTok{ p}\OperatorTok{.}\AttributeTok{coords}\OperatorTok{.}\AttributeTok{latitude}\OperatorTok{,} \DataTypeTok{lng}\OperatorTok{:}\NormalTok{ p}\OperatorTok{.}\AttributeTok{coords}\OperatorTok{.}\AttributeTok{longitude}\NormalTok{ \})}\OperatorTok{,}
\NormalTok{      ()   }\KeywordTok{=\textgreater{}} \FunctionTok{resolve}\NormalTok{(fallback)}\OperatorTok{,}
\NormalTok{      \{ }\DataTypeTok{timeout}\OperatorTok{:} \DecValTok{8000}\OperatorTok{,} \DataTypeTok{maximumAge}\OperatorTok{:} \DecValTok{600000}\NormalTok{ \})}\OperatorTok{;}
\NormalTok{  \})}\OperatorTok{;}
\NormalTok{\}}

\KeywordTok{async} \KeywordTok{function} \FunctionTok{connect}\NormalTok{() \{}
  \KeywordTok{const}\NormalTok{ here }\OperatorTok{=} \ControlFlowTok{await} \FunctionTok{whereAmI}\NormalTok{()}\OperatorTok{;}

  \CommentTok{// CONNECTION identity {-}{-} ephemeral, rooted at the user\textquotesingle{}s real location.}
\NormalTok{  nodeIdentity }\OperatorTok{=} \ControlFlowTok{await} \FunctionTok{createNodeIdentity}\NormalTok{(\{ }\DataTypeTok{lat}\OperatorTok{:}\NormalTok{ here}\OperatorTok{.}\AttributeTok{lat}\OperatorTok{,} \DataTypeTok{lng}\OperatorTok{:}\NormalTok{ here}\OperatorTok{.}\AttributeTok{lng}\NormalTok{ \})}\OperatorTok{;}

  \CommentTok{// AUTHORSHIP identity {-}{-} durable, location{-}free, names the signer on every post.}
\NormalTok{  author }\OperatorTok{=} \ControlFlowTok{await} \FunctionTok{createAuthorIdentity}\NormalTok{(\{ }\DataTypeTok{persistAs}\OperatorTok{:} \StringTok{\textquotesingle{}axonatalk:author\textquotesingle{}}\NormalTok{ \})}\OperatorTok{;}

  \KeywordTok{const}\NormalTok{ transport }\OperatorTok{=} \FunctionTok{webTransport}\NormalTok{(\{ }\DataTypeTok{bridgeUrl}\OperatorTok{:}\NormalTok{ BRIDGE}\OperatorTok{,} \DataTypeTok{identity}\OperatorTok{:}\NormalTok{ nodeIdentity \})}\OperatorTok{;}
  \KeywordTok{const}\NormalTok{ node }\OperatorTok{=} \KeywordTok{new} \FunctionTok{NeuronNode}\NormalTok{(\{ }\DataTypeTok{id}\OperatorTok{:} \BuiltInTok{BigInt}\NormalTok{(}\StringTok{\textquotesingle{}0x\textquotesingle{}} \OperatorTok{+}\NormalTok{ nodeIdentity}\OperatorTok{.}\AttributeTok{id}\NormalTok{)}\OperatorTok{,} \DataTypeTok{lat}\OperatorTok{:}\NormalTok{ here}\OperatorTok{.}\AttributeTok{lat}\OperatorTok{,} \DataTypeTok{lng}\OperatorTok{:}\NormalTok{ here}\OperatorTok{.}\AttributeTok{lng}\NormalTok{ \})}\OperatorTok{;}
\NormalTok{  node}\OperatorTok{.}\AttributeTok{transport} \OperatorTok{=}\NormalTok{ transport}\OperatorTok{;}
\NormalTok{  peer }\OperatorTok{=} \KeywordTok{new} \FunctionTok{AxonaPeer}\NormalTok{(\{ }\DataTypeTok{domain}\OperatorTok{:} \KeywordTok{new} \FunctionTok{AxonaDomain}\NormalTok{(\{ }\DataTypeTok{k}\OperatorTok{:} \DecValTok{20}\NormalTok{ \})}\OperatorTok{,}\NormalTok{ node}\OperatorTok{,}\NormalTok{ nodeIdentity}\OperatorTok{,}\NormalTok{ transport \})}\OperatorTok{;}

  \ControlFlowTok{await}\NormalTok{ transport}\OperatorTok{.}\FunctionTok{start}\NormalTok{(nodeIdentity}\OperatorTok{.}\AttributeTok{id}\NormalTok{)}\OperatorTok{;}
  \ControlFlowTok{await}\NormalTok{ peer}\OperatorTok{.}\FunctionTok{start}\NormalTok{()}\OperatorTok{;}

  \CommentTok{// Wait for the mesh to form before publishing (see pitfall 13.3). peer.ready()}
  \CommentTok{// does this for you {-}{-} it resolves once enough peers are in the synaptome (or a}
  \CommentTok{// stable non{-}zero plateau), bounded by timeoutMs. No hand{-}rolled polling loop.}
  \KeywordTok{const}\NormalTok{ status }\OperatorTok{=} \ControlFlowTok{await}\NormalTok{ peer}\OperatorTok{.}\FunctionTok{ready}\NormalTok{()}\OperatorTok{;}   \CommentTok{// \{ ready, peers, ms, reason \}}
  \BuiltInTok{console}\OperatorTok{.}\FunctionTok{log}\NormalTok{(}\VerbatimStringTok{\textasciigrave{}AxonaTalk on kernel }\SpecialCharTok{$\{}\NormalTok{KERNEL\_VERSION}\SpecialCharTok{\}}\VerbatimStringTok{; you are }\SpecialCharTok{$\{}\FunctionTok{idLabel}\NormalTok{(nodeIdentity}\OperatorTok{.}\AttributeTok{id}\NormalTok{)}\SpecialCharTok{\}}\VerbatimStringTok{ \textasciigrave{}} \OperatorTok{+}
              \VerbatimStringTok{\textasciigrave{}(mesh }\SpecialCharTok{$\{}\NormalTok{status}\OperatorTok{.}\AttributeTok{ready} \OperatorTok{?} \StringTok{\textquotesingle{}ready\textquotesingle{}} \OperatorTok{:} \StringTok{\textquotesingle{}not ready\textquotesingle{}}\SpecialCharTok{\}}\VerbatimStringTok{, }\SpecialCharTok{$\{}\NormalTok{status}\OperatorTok{.}\AttributeTok{peers}\SpecialCharTok{\}}\VerbatimStringTok{ peers, }\SpecialCharTok{$\{}\NormalTok{status}\OperatorTok{.}\AttributeTok{ms}\SpecialCharTok{\}}\VerbatimStringTok{ms)\textasciigrave{}}\NormalTok{)}\OperatorTok{;}
  \ControlFlowTok{return}\NormalTok{ \{ peer}\OperatorTok{,}\NormalTok{ author}\OperatorTok{,}\NormalTok{ nodeIdentity \}}\OperatorTok{;}
\NormalTok{\}}

\CommentTok{// "region:idprefix" read off a 264{-}bit Node ID we VOLUNTARILY shared in the}
\CommentTok{// payload {-}{-} top byte = the sender\textquotesingle{}s S2 region. The protocol itself never}
\CommentTok{// carries this. Anonymous / absent {-}\textgreater{} \textquotesingle{}anon\textquotesingle{}.}
\KeywordTok{const}\NormalTok{ idLabel }\OperatorTok{=}\NormalTok{ (nodeId) }\KeywordTok{=\textgreater{}}
\NormalTok{  (}\KeywordTok{typeof}\NormalTok{ nodeId }\OperatorTok{===} \StringTok{\textquotesingle{}string\textquotesingle{}} \OperatorTok{\&\&}\NormalTok{ nodeId}\OperatorTok{.}\AttributeTok{length} \OperatorTok{\textgreater{}=} \DecValTok{10}\NormalTok{)}
    \OperatorTok{?} \VerbatimStringTok{\textasciigrave{}}\SpecialCharTok{$\{}\FunctionTok{regionName}\NormalTok{(}\PreprocessorTok{parseInt}\NormalTok{(nodeId}\OperatorTok{.}\FunctionTok{slice}\NormalTok{(}\DecValTok{0}\OperatorTok{,} \DecValTok{2}\NormalTok{)}\OperatorTok{,} \DecValTok{16}\NormalTok{))}\SpecialCharTok{\}}\VerbatimStringTok{:}\SpecialCharTok{$\{}\NormalTok{nodeId}\OperatorTok{.}\FunctionTok{slice}\NormalTok{(}\DecValTok{2}\OperatorTok{,} \DecValTok{10}\NormalTok{)}\SpecialCharTok{\}}\VerbatimStringTok{\textasciigrave{}}
    \OperatorTok{:} \StringTok{\textquotesingle{}anon\textquotesingle{}}\OperatorTok{;}
\end{Highlighting}
\end{Shaded}

\subsubsection{12.4 Join a room}\label{join-a-room}

\begin{Shaded}
\begin{Highlighting}[]
\KeywordTok{let}\NormalTok{ currentRoom }\OperatorTok{=} \KeywordTok{null}\OperatorTok{,}\NormalTok{ currentSub }\OperatorTok{=} \KeywordTok{null}\OperatorTok{;}
\KeywordTok{const}\NormalTok{ seen }\OperatorTok{=} \KeywordTok{new} \BuiltInTok{Set}\NormalTok{()}\OperatorTok{;}   \CommentTok{// msgIds already shown (dedup our own echo)}

\KeywordTok{async} \KeywordTok{function} \FunctionTok{joinRoom}\NormalTok{(room}\OperatorTok{,}\NormalTok{ onMessage) \{}
  \ControlFlowTok{if}\NormalTok{ (room }\OperatorTok{===}\NormalTok{ currentRoom) }\ControlFlowTok{return}\OperatorTok{;}
  \ControlFlowTok{if}\NormalTok{ (currentSub) \{ }\ControlFlowTok{try}\NormalTok{ \{ }\ControlFlowTok{await}\NormalTok{ currentSub}\OperatorTok{.}\FunctionTok{stop}\NormalTok{()}\OperatorTok{;}\NormalTok{ \} }\ControlFlowTok{catch}\NormalTok{ \{\} \}}
\NormalTok{  currentRoom }\OperatorTok{=}\NormalTok{ room}\OperatorTok{;}

\NormalTok{  currentSub }\OperatorTok{=} \ControlFlowTok{await}\NormalTok{ peer}\OperatorTok{.}\FunctionTok{sub}\NormalTok{(\{ }\DataTypeTok{region}\OperatorTok{:}\NormalTok{ TOPIC\_REGION}\OperatorTok{,} \DataTypeTok{name}\OperatorTok{:}\NormalTok{ room \}}\OperatorTok{,}\NormalTok{ (env) }\KeywordTok{=\textgreater{}}\NormalTok{ \{}
    \ControlFlowTok{if}\NormalTok{ (}\OperatorTok{!}\NormalTok{env }\OperatorTok{||}\NormalTok{ env}\OperatorTok{.}\AttributeTok{deleted} \OperatorTok{||}\NormalTok{ seen}\OperatorTok{.}\FunctionTok{has}\NormalTok{(env}\OperatorTok{.}\AttributeTok{msgId}\NormalTok{)) }\ControlFlowTok{return}\OperatorTok{;}
\NormalTok{    seen}\OperatorTok{.}\FunctionTok{add}\NormalTok{(env}\OperatorTok{.}\AttributeTok{msgId}\NormalTok{)}\OperatorTok{;}
    \KeywordTok{const}\NormalTok{ m }\OperatorTok{=}\NormalTok{ env}\OperatorTok{.}\AttributeTok{message}\OperatorTok{;}
    \FunctionTok{onMessage}\NormalTok{(\{}
      \DataTypeTok{text}\OperatorTok{:}\NormalTok{ (m }\OperatorTok{\&\&} \KeywordTok{typeof}\NormalTok{ m }\OperatorTok{===} \StringTok{\textquotesingle{}object\textquotesingle{}}\NormalTok{) }\OperatorTok{?}\NormalTok{ m}\OperatorTok{.}\AttributeTok{text} \OperatorTok{:}\NormalTok{ m}\OperatorTok{,}
      \DataTypeTok{who}\OperatorTok{:}  \FunctionTok{idLabel}\NormalTok{((m }\OperatorTok{\&\&} \KeywordTok{typeof}\NormalTok{ m }\OperatorTok{===} \StringTok{\textquotesingle{}object\textquotesingle{}}\NormalTok{) }\OperatorTok{?}\NormalTok{ m}\OperatorTok{.}\AttributeTok{node} \OperatorTok{:} \KeywordTok{null}\NormalTok{)}\OperatorTok{,}
\NormalTok{    \})}\OperatorTok{;}
\NormalTok{  \}}\OperatorTok{,}\NormalTok{ \{ }\DataTypeTok{since}\OperatorTok{:} \StringTok{\textquotesingle{}all\textquotesingle{}}\NormalTok{ \})}\OperatorTok{;}   \CommentTok{// give new joiners the room history}
\NormalTok{\}}
\end{Highlighting}
\end{Shaded}

\subsubsection{12.5 Send a chat message}\label{send-a-chat-message}

\begin{Shaded}
\begin{Highlighting}[]
\KeywordTok{async} \KeywordTok{function} \FunctionTok{sendChat}\NormalTok{(room}\OperatorTok{,}\NormalTok{ text) \{}
  \ControlFlowTok{await} \FunctionTok{joinRoom}\NormalTok{(room}\OperatorTok{,}\NormalTok{ renderIncoming)}\OperatorTok{;}   \CommentTok{// make sure we\textquotesingle{}re subscribed}
  \CommentTok{// Every publish names its signer. We voluntarily include our Node ID so}
  \CommentTok{// subscribers can show our region; the signed envelope never carries it.}
  \KeywordTok{const}\NormalTok{ msgId }\OperatorTok{=} \ControlFlowTok{await}\NormalTok{ peer}\OperatorTok{.}\FunctionTok{pub}\NormalTok{(}
\NormalTok{    \{ }\DataTypeTok{region}\OperatorTok{:}\NormalTok{ TOPIC\_REGION}\OperatorTok{,} \DataTypeTok{name}\OperatorTok{:}\NormalTok{ room \}}\OperatorTok{,}
\NormalTok{    \{ text}\OperatorTok{,} \DataTypeTok{node}\OperatorTok{:}\NormalTok{ nodeIdentity}\OperatorTok{.}\AttributeTok{id}\NormalTok{ \}}\OperatorTok{,}
\NormalTok{    \{ }\DataTypeTok{signWith}\OperatorTok{:}\NormalTok{ author \})}\OperatorTok{;}
\NormalTok{  seen}\OperatorTok{.}\FunctionTok{add}\NormalTok{(msgId)}\OperatorTok{;}                        \CommentTok{// our own publish may not echo back}
  \FunctionTok{renderIncoming}\NormalTok{(\{ text}\OperatorTok{,} \DataTypeTok{who}\OperatorTok{:} \FunctionTok{idLabel}\NormalTok{(nodeIdentity}\OperatorTok{.}\AttributeTok{id}\NormalTok{)}\OperatorTok{,} \DataTypeTok{self}\OperatorTok{:} \KeywordTok{true}\NormalTok{ \})}\OperatorTok{;}
  \ControlFlowTok{return}\NormalTok{ msgId}\OperatorTok{;}
\NormalTok{\}}
\end{Highlighting}
\end{Shaded}

\subsubsection{12.6 Show a room's Topic ID (shareable read
handle)}\label{show-a-rooms-topic-id-shareable-read-handle}

\begin{Shaded}
\begin{Highlighting}[]
\KeywordTok{async} \KeywordTok{function} \FunctionTok{roomLink}\NormalTok{(room) \{}
  \CommentTok{// deriveTopicId is pure {-}{-} works before we even connect.}
  \KeywordTok{const}\NormalTok{ id }\OperatorTok{=} \ControlFlowTok{await} \FunctionTok{deriveTopicId}\NormalTok{(\{ }\DataTypeTok{region}\OperatorTok{:}\NormalTok{ TOPIC\_REGION}\OperatorTok{,} \DataTypeTok{name}\OperatorTok{:}\NormalTok{ room \})}\OperatorTok{;}
  \ControlFlowTok{return} \VerbatimStringTok{\textasciigrave{}}\SpecialCharTok{$\{}\FunctionTok{regionName}\NormalTok{(}\BaseNTok{0x89}\NormalTok{)}\SpecialCharTok{\}}\VerbatimStringTok{ : }\SpecialCharTok{$\{}\NormalTok{id}\SpecialCharTok{\}}\VerbatimStringTok{\textasciigrave{}}\OperatorTok{;}   \CommentTok{// share this so others can sub by ID}
\NormalTok{\}}
\end{Highlighting}
\end{Shaded}

\subsubsection{12.7 Direct messages}\label{direct-messages}

\begin{Shaded}
\begin{Highlighting}[]
\KeywordTok{async} \KeywordTok{function} \FunctionTok{sendDM}\NormalTok{(recipientNodeIdHex}\OperatorTok{,}\NormalTok{ text) \{}
  \ControlFlowTok{return}\NormalTok{ peer}\OperatorTok{.}\FunctionTok{send}\NormalTok{(recipientNodeIdHex}\OperatorTok{,}\NormalTok{ \{}
    \DataTypeTok{kind}\OperatorTok{:} \StringTok{\textquotesingle{}axonatalk:dm\textquotesingle{}}\OperatorTok{,} \DataTypeTok{from}\OperatorTok{:}\NormalTok{ author}\OperatorTok{.}\AttributeTok{authorId}\OperatorTok{,}\NormalTok{ text}\OperatorTok{,} \DataTypeTok{ts}\OperatorTok{:} \BuiltInTok{Date}\OperatorTok{.}\FunctionTok{now}\NormalTok{()}\OperatorTok{,}
\NormalTok{  \})}\OperatorTok{;}
\NormalTok{\}}

\KeywordTok{function} \FunctionTok{registerDMHandler}\NormalTok{(onDM) \{}
\NormalTok{  peer}\OperatorTok{.}\FunctionTok{onMessage}\NormalTok{(}\KeywordTok{async}\NormalTok{ (senderId}\OperatorTok{,}\NormalTok{ message) }\KeywordTok{=\textgreater{}}\NormalTok{ \{}
    \ControlFlowTok{if}\NormalTok{ (message}\OperatorTok{?.}\AttributeTok{kind} \OperatorTok{===} \StringTok{\textquotesingle{}axonatalk:dm\textquotesingle{}}\NormalTok{) \{}
      \FunctionTok{onDM}\NormalTok{(\{ }\DataTypeTok{fromNode}\OperatorTok{:}\NormalTok{ senderId}\OperatorTok{,} \DataTypeTok{fromAuthor}\OperatorTok{:}\NormalTok{ message}\OperatorTok{.}\AttributeTok{from}\OperatorTok{,} \DataTypeTok{text}\OperatorTok{:}\NormalTok{ message}\OperatorTok{.}\AttributeTok{text}\NormalTok{ \})}\OperatorTok{;}
      \ControlFlowTok{return}\NormalTok{ \{ }\DataTypeTok{ack}\OperatorTok{:} \KeywordTok{true}\OperatorTok{,} \DataTypeTok{receivedAt}\OperatorTok{:} \BuiltInTok{Date}\OperatorTok{.}\FunctionTok{now}\NormalTok{() \}}\OperatorTok{;}
\NormalTok{    \}}
\NormalTok{  \})}\OperatorTok{;}
\NormalTok{\}}
\end{Highlighting}
\end{Shaded}

\subsubsection{12.8 What you did and did not have to
build}\label{what-you-did-and-did-not-have-to-build}

You did \textbf{not} build: a server (the bridge is shared infra),
replay/history
(\texttt{since:\ \textquotesingle{}all\textquotesingle{}}), a message
bus (the synaptome routes), or cryptography (\texttt{signWith} signs,
\texttt{verifyEnvelope} checks).

You \textbf{did} decide: what an author means in your social model, your
trust policy (who you accept by Author ID), and the UI. None of that is
the protocol's job.

\begin{center}\rule{0.5\linewidth}{0.5pt}\end{center}

\subsection{13. Common pitfalls}\label{common-pitfalls}

\subsubsection{\texorpdfstring{13.1 Forgetting
\texttt{signWith}}{13.1 Forgetting signWith}}\label{forgetting-signwith}

\texttt{peer.pub} has \textbf{no default signer.} Omitting
\texttt{signWith} throws \texttt{PUBLISH\_NO\_PUBLISH\_IDENTITY} -- it
does not silently publish anonymously. Pass
\texttt{\{\ signWith:\ author\ \}}, or
\texttt{\{\ signWith:\ ANONYMOUS\ \}} if you really mean unsigned. (This
trips up everyone migrating from the v2 line, where the node key signed
by default.)

\subsubsection{13.2 Publishing with a Topic
ID}\label{publishing-with-a-topic-id}

\texttt{pub} and \texttt{kill} need the \textbf{descriptor} -- a bare
66-hex Topic ID throws \texttt{PUBLISH\_INVALID\_TOPIC}. Only
\texttt{sub}, \texttt{pull}, and \texttt{metrics} accept an ID. Share
the ID for reading; share the descriptor (\texttt{\{\ region,\ name\ \}}
or the full owned descriptor) for writing (6.6).

\subsubsection{13.3 Publishing before the mesh
forms}\label{publishing-before-the-mesh-forms}

\texttt{peer.start()} returns as soon as local state is set up; mesh
handshakes complete asynchronously over the next few seconds. Publish
too early and your K-closest set is just yourself + the bridge, so the
axon set is incomplete. The 4.x \textbf{cold-publish burst} (§7.1) now
re-sends a freshly-joined node's first publishes over the first second,
so an early publish usually still lands --- this is a soft edge, not the
hard 0\% it once was. Still, for predictable UX,
\texttt{await\ peer.ready()} before letting the user publish --- it
resolves once the synaptome reaches a small threshold (or a stable
non-zero plateau), bounded by \texttt{timeoutMs}, and hands back
\texttt{\{\ ready,\ peers,\ ms,\ reason\ \}}. The worked example (12.3)
uses it; prefer it over hand-rolling a \texttt{node.synaptome.size} /
\texttt{peer.peers().length} poll.

\subsubsection{13.4 Mismatched descriptors give no delivery,
silently}\label{mismatched-descriptors-give-no-delivery-silently}

Publisher and subscriber must resolve to the \textbf{same Topic ID.} If
one names \texttt{region:\ \textquotesingle{}useast\textquotesingle{}}
and the other omits \texttt{region} (defaulting to a different node
region), or one adds an \texttt{owner} the other does not, the IDs
differ and nothing is delivered -- with no warning. Centralize your
descriptor construction in one helper and call it from both sides.

\subsubsection{13.5 Owned-topic write
policy}\label{owned-topic-write-policy}

Publishing to \texttt{\{\ owner,\ name\ \}} (write defaults to
\texttt{owner}) with a \texttt{signWith} that is not the owner key
throws \texttt{WRITE\_POLICY\_VIOLATION} (the peer fails fast; the
storing node enforces the same at ingress). If you meant an inbox anyone
can post to, set
\texttt{write:\ \textquotesingle{}open\textquotesingle{}} explicitly.

\subsubsection{13.6 Replay needs the message to still be cached -- and
someone to host
it}\label{replay-needs-the-message-to-still-be-cached-and-someone-to-host-it}

\texttt{since:\ \textquotesingle{}all\textquotesingle{}} returns
whatever is in the axons' caches \emph{right now}. If a topic has no
subscribers and no hosting node, its role gets swept after a short grace
window and the cache vanishes -- even before the hold time elapses. For
durable history, run a hosting node: \texttt{peer.host()} on a
long-lived peer (or use \texttt{axona-relay}), or mirror envelopes into
your own database at the app layer.

\subsubsection{13.7 Message size}\label{message-size}

The reliable-publish floor is \textbf{16 KiB} on the serialized
envelope. \texttt{peer.pub} throws \texttt{PUBLISH\_PAYLOAD\_TOO\_LARGE}
above it rather than letting an unreceivable message vanish mid-mesh. Do
not ship blobs over pub/sub even under the cap -- every publish
replicates to K axons and fans to all subscribers. Publish a small
content reference and move the bytes out of band, or use
\texttt{@axona/protocol/std/chunk}.

\subsubsection{13.8 Node identity is ephemeral; persist the author, not
the
node}\label{node-identity-is-ephemeral-persist-the-author-not-the-node}

A fresh \texttt{createNodeIdentity} each run means a new Node ID each
run -- by design (4.2). If your app needs stable identity for ``block
this user'' lists or persistent threads, that stability comes from the
\textbf{author} key (\texttt{persistAs}), which is location-free and is
the recognizable thing. Persist the node identity only if you
specifically want a stable address (rare).

\subsubsection{\texorpdfstring{13.9 \texttt{verifyEnvelope} returns an
object}{13.9 verifyEnvelope returns an object}}\label{verifyenvelope-returns-an-object}

Test \texttt{.ok}. \texttt{if\ (await\ verifyEnvelope(env))} is always
truthy and would never reject a forgery (5.5).

\subsubsection{13.10 Hard-reload after a
deploy}\label{hard-reload-after-a-deploy}

Browser bundle cache is sticky. After deploying a new peer, hard-reload
(Cmd+Shift+R / Ctrl+Shift+R). If tabs run different versions across a
wire change, the mesh half-works and debugging is miserable. The bridge
rejects peers below its minimum version with an upgrade close code,
surfaced via \texttt{peer.onUpgradeRequired}.

\begin{center}\rule{0.5\linewidth}{0.5pt}\end{center}

\subsection{14. Production checklist}\label{production-checklist}

Before launching a public Axona app:

\begin{itemize}
\tightlist
\item[$\square$]
  \textbf{Docker bridge stack} --
  \texttt{docker\ compose\ up\ -d\ -\/-build} with \texttt{DOMAIN},
  \texttt{PUBLIC\_IP}, \texttt{TURN\_AUTH\_SECRET},
  \texttt{BRIDGE\_PUBLIC\_URL} set. Caddy gives you automatic HTTPS; do
  not hand-roll TLS.
\item[$\square$]
  \textbf{TURN alongside the bridge} -- coturn with a shared
  \texttt{TURN\_AUTH\_SECRET}. Not optional: symmetric-NAT peers cannot
  connect without it.
\item[$\square$]
  \textbf{Bridge directory / federation} -- set
  \texttt{BRIDGE\_PUBLIC\_URL} so the bridge advertises itself; set
  \texttt{BRIDGE\_DIRECTORY=off} for a private fleet. Verify
  \texttt{/healthz} shows the directory and uplink state you expect.
\item[$\square$]
  \textbf{Stable bridge identity} -- the persistent \texttt{bridge-data}
  volume keeps the bridge's Node ID across restarts.
\item[$\square$]
  \textbf{Author persistence} --
  \texttt{createAuthorIdentity(\{\ persistAs\ \})} so users keep a
  recognizable authorship across sessions.
\item[$\square$]
  \textbf{Descriptor-based topic naming} -- decide and document your
  topic shapes: which are open lobbies \texttt{\{\ region,\ name\ \}},
  which are owned feeds \texttt{\{\ region,\ owner,\ name\ \}}, which
  are inboxes
  (\texttt{write:\ \textquotesingle{}open\textquotesingle{}}). Build
  every descriptor through one shared helper so pub and sub never drift
  (13.4).
\item[$\square$]
  \textbf{Region policy} -- decide whether topics pin to one named
  region (for a shared keyspace) or default to the node region, and
  apply it consistently.
\item[$\square$]
  \textbf{Signature / trust policy} -- decide which Author IDs you trust
  and reject the rest in your handler (or accept everyone). Verify with
  \texttt{verifyEnvelope}.
\item[$\square$]
  \textbf{Durable history} -- if you need history beyond the cache
  window, run a hosting node (\texttt{peer.host()} /
  \texttt{axona-relay}) or journal envelopes to your own store (13.6).
\item[$\square$]
  \textbf{Surface \texttt{KERNEL\_VERSION}} -- in the app version row
  and at the bridge \texttt{/healthz}, imported from the kernel (11.5).
\item[$\square$]
  \textbf{Health monitoring} -- scrape
  \texttt{https://\textless{}bridge\textgreater{}/healthz} and/or hook
  \texttt{peer.health().meshDegraded} into your dashboards.
\end{itemize}

\begin{center}\rule{0.5\linewidth}{0.5pt}\end{center}

\subsection{15. Cheat sheet}\label{cheat-sheet}

All imports are
\texttt{from\ \textquotesingle{}@axona/protocol\textquotesingle{}}
unless a sub-path is shown.

\subsubsection{15.1 Identity}\label{identity}

\begin{Shaded}
\begin{Highlighting}[]
\FunctionTok{createNodeIdentity}\NormalTok{(\{ lat}\OperatorTok{,}\NormalTok{ lng}\OperatorTok{,}\NormalTok{ extractable}\OperatorTok{?}\NormalTok{ \})   }\CommentTok{// {-}\textgreater{} connection identity (Node ID)}
\FunctionTok{createAuthorIdentity}\NormalTok{(\{ persistAs}\OperatorTok{?,}\NormalTok{ store}\OperatorTok{?,}\NormalTok{ extractable}\OperatorTok{?}\NormalTok{ \})  }\CommentTok{// {-}\textgreater{} author identity (Author ID)}
\NormalTok{ANONYMOUS                                          }\CommentTok{// sentinel for an unsigned publish}

\FunctionTok{dumpIdentity}\NormalTok{(nodeIdentity)  }\OperatorTok{/} \FunctionTok{loadIdentity}\NormalTok{(env)    }\CommentTok{// node{-}identity persistence envelope}
\end{Highlighting}
\end{Shaded}

\subsubsection{15.2 Peer}\label{peer}

\begin{Shaded}
\begin{Highlighting}[]
\KeywordTok{new} \FunctionTok{AxonaPeer}\NormalTok{(\{ domain}\OperatorTok{,}\NormalTok{ node}\OperatorTok{,}\NormalTok{ nodeIdentity}\OperatorTok{,}\NormalTok{ transport}\OperatorTok{,}\NormalTok{ persist}\OperatorTok{?}\NormalTok{ \})   }\CommentTok{// NO author key here}
\ControlFlowTok{await}\NormalTok{ peer}\OperatorTok{.}\FunctionTok{start}\NormalTok{()}\OperatorTok{;}  \ControlFlowTok{await}\NormalTok{ peer}\OperatorTok{.}\FunctionTok{stop}\NormalTok{()}\OperatorTok{;}
\ControlFlowTok{await}\NormalTok{ peer}\OperatorTok{.}\FunctionTok{join}\NormalTok{(sponsorHex}\OperatorTok{?}\NormalTok{)}\OperatorTok{;}  \ControlFlowTok{await}\NormalTok{ peer}\OperatorTok{.}\FunctionTok{leave}\NormalTok{(\{ drain}\OperatorTok{?,}\NormalTok{ notify}\OperatorTok{?,}\NormalTok{ timeoutMs}\OperatorTok{?}\NormalTok{ \})}\OperatorTok{;}
\NormalTok{peer}\OperatorTok{.}\FunctionTok{getNodeId}\NormalTok{()}\OperatorTok{;}   \CommentTok{// 66{-}hex Node ID}
\end{Highlighting}
\end{Shaded}

\subsubsection{15.3 Topics}\label{topics}

\begin{Shaded}
\begin{Highlighting}[]
\FunctionTok{deriveTopicId}\NormalTok{(\{ region}\OperatorTok{?,}\NormalTok{ owner}\OperatorTok{?,}\NormalTok{ name}\OperatorTok{,}\NormalTok{ write}\OperatorTok{?}\NormalTok{ \})   }\CommentTok{// {-}\textgreater{} 66{-}hex read handle}
\CommentTok{// descriptor: region = name | \textquotesingle{}0x89\textquotesingle{} | 137 | omit(node region); owner = Author ID;}
\CommentTok{//             write = \textquotesingle{}open\textquotesingle{} | \textquotesingle{}owner\textquotesingle{} (defaults: owner present {-}\textgreater{} \textquotesingle{}owner\textquotesingle{}, else \textquotesingle{}open\textquotesingle{})}
\end{Highlighting}
\end{Shaded}

\subsubsection{15.4 Pub/Sub}\label{pubsub}

\begin{Shaded}
\begin{Highlighting}[]
\NormalTok{peer}\OperatorTok{.}\FunctionTok{pub}\NormalTok{(descriptor}\OperatorTok{,}\NormalTok{ message}\OperatorTok{,}\NormalTok{ \{ signWith \})        }\CommentTok{// {-}\textgreater{} msgId   (signWith REQUIRED)}
\NormalTok{peer}\OperatorTok{.}\FunctionTok{sub}\NormalTok{(descriptor }\OperatorTok{|}\NormalTok{ topicId}\OperatorTok{,}\NormalTok{ handler}\OperatorTok{,}\NormalTok{ \{ since}\OperatorTok{?}\NormalTok{ \}) }\CommentTok{// {-}\textgreater{} Subscription; handler gets envelope or \{ msgId, deleted:true \}}
\NormalTok{peer}\OperatorTok{.}\FunctionTok{unsub}\NormalTok{(descriptor)                              }\CommentTok{// {-}\textgreater{} \{ ok, removed \}}
\NormalTok{peer}\OperatorTok{.}\FunctionTok{kill}\NormalTok{(descriptor}\OperatorTok{,}\NormalTok{ msgId}\OperatorTok{,}\NormalTok{ \{ signWith \})          }\CommentTok{// author{-}only retract (sole retraction primitive; touch/unpub removed v4.3.0)}
\NormalTok{peer}\OperatorTok{.}\FunctionTok{pull}\NormalTok{(msgId }\OperatorTok{|} \KeywordTok{null}\OperatorTok{,}\NormalTok{ \{ topic}\OperatorTok{,}\NormalTok{ timeoutMs}\OperatorTok{?}\NormalTok{ \})      }\CommentTok{// topic = descriptor | id; null {-}\textgreater{} latest}
\NormalTok{peer}\OperatorTok{.}\FunctionTok{metrics}\NormalTok{(descriptor }\OperatorTok{|}\NormalTok{ topicId}\OperatorTok{,}\NormalTok{ \{ timeoutMs}\OperatorTok{?}\NormalTok{ \})  }\CommentTok{// one{-}shot read of the published snapshot (open + owned)}
\FunctionTok{metricTopic}\NormalTok{(dataTopicId)                            }\CommentTok{// {-}\textgreater{} metric topic descriptor; sub() it for live + trend}
\NormalTok{peer}\OperatorTok{.}\FunctionTok{sub}\NormalTok{(}\FunctionTok{metricTopic}\NormalTok{(}\ControlFlowTok{await} \FunctionTok{deriveTopicId}\NormalTok{(desc))}\OperatorTok{,}\NormalTok{ cb}\OperatorTok{,}\NormalTok{ \{ }\DataTypeTok{since}\OperatorTok{:}\StringTok{\textquotesingle{}all\textquotesingle{}}\NormalTok{ \})   }\CommentTok{// scalable metrics}
\NormalTok{peer}\OperatorTok{.}\FunctionTok{host}\NormalTok{(descriptor}\OperatorTok{?}\NormalTok{)  }\OperatorTok{/}\NormalTok{ peer}\OperatorTok{.}\FunctionTok{unhost}\NormalTok{(descriptor}\OperatorTok{?}\NormalTok{)  }\CommentTok{// host()/unhost() = keyspace}
\NormalTok{peer}\OperatorTok{.}\FunctionTok{rootedTopics}\NormalTok{()                                 }\CommentTok{// infra: topics I root + local snapshots}
\ControlFlowTok{await}\NormalTok{ subscription}\OperatorTok{.}\FunctionTok{stop}\NormalTok{()}\OperatorTok{;}
\CommentTok{// since: omit (live) | \textquotesingle{}all\textquotesingle{} | \textquotesingle{}latest\textquotesingle{} | \textless{}ms epoch\textgreater{}}
\end{Highlighting}
\end{Shaded}

\subsubsection{15.5 Direct messaging}\label{direct-messaging-1}

\begin{Shaded}
\begin{Highlighting}[]
\NormalTok{peer}\OperatorTok{.}\FunctionTok{send}\NormalTok{(nodeIdHex}\OperatorTok{,}\NormalTok{ message)    }\CommentTok{// {-}\textgreater{} remote handler\textquotesingle{}s return value}
\NormalTok{peer}\OperatorTok{.}\FunctionTok{notify}\NormalTok{(nodeIdHex}\OperatorTok{,}\NormalTok{ message)  }\CommentTok{// fire{-}and{-}forget}
\NormalTok{peer}\OperatorTok{.}\FunctionTok{onMessage}\NormalTok{((senderId}\OperatorTok{,}\NormalTok{ message) }\KeywordTok{=\textgreater{}}\NormalTok{ reply)   }\CommentTok{// single handler}
\end{Highlighting}
\end{Shaded}

\subsubsection{15.6 Introspection}\label{introspection}

\begin{Shaded}
\begin{Highlighting}[]
\NormalTok{peer}\OperatorTok{.}\FunctionTok{peers}\NormalTok{()}\OperatorTok{;}                    \CommentTok{// {-}\textgreater{} Node ID[] in the synaptome}
\NormalTok{peer}\OperatorTok{.}\FunctionTok{onPeerJoin}\NormalTok{(h)}\OperatorTok{;}\NormalTok{  peer}\OperatorTok{.}\FunctionTok{onPeerLeave}\NormalTok{(h)}\OperatorTok{;}   \CommentTok{// {-}\textgreater{} unsubscribe()}
\NormalTok{peer}\OperatorTok{.}\FunctionTok{lookup}\NormalTok{(targetBigInt)}\OperatorTok{;}       \CommentTok{// {-}\textgreater{} \{ found, hops, time, path \}}
\NormalTok{peer}\OperatorTok{.}\FunctionTok{health}\NormalTok{()}\OperatorTok{;}                   \CommentTok{// {-}\textgreater{} diagnostic snapshot (meshDegraded, transport, axonRoles, ...)}
\NormalTok{peer}\OperatorTok{.}\FunctionTok{onLog}\NormalTok{(level}\OperatorTok{,}\NormalTok{ h)}\OperatorTok{;}\NormalTok{  peer}\OperatorTok{.}\FunctionTok{onError}\NormalTok{(h)}\OperatorTok{;}\NormalTok{  peer}\OperatorTok{.}\FunctionTok{onUpgradeRequired}\NormalTok{(h)}\OperatorTok{;}
\end{Highlighting}
\end{Shaded}

\subsubsection{15.7 Envelope + region
helpers}\label{envelope-region-helpers}

\begin{Shaded}
\begin{Highlighting}[]
\FunctionTok{verifyEnvelope}\NormalTok{(envelope)         }\CommentTok{// {-}\textgreater{} \{ ok, reason?, signed \}  (test .ok, not truthiness)}
\FunctionTok{regionName}\NormalTok{(code)}\OperatorTok{;}  \FunctionTok{regionCode}\NormalTok{(name)}\OperatorTok{;}  \FunctionTok{resolveRegion}\NormalTok{(token)}\OperatorTok{;}  \FunctionTok{regionCenter}\NormalTok{(nameOrCode)}\OperatorTok{;}
\end{Highlighting}
\end{Shaded}

\subsubsection{15.8 Persistence (sub-path
imports)}\label{persistence-sub-path-imports}

\begin{Shaded}
\begin{Highlighting}[]
\ImportTok{import}\NormalTok{ \{ IndexedDBPersistence \} }\ImportTok{from} \StringTok{\textquotesingle{}@axona/protocol/persistence/indexeddb.js\textquotesingle{}}\OperatorTok{;}
\ImportTok{import}\NormalTok{ \{ FilePersistence \}      }\ImportTok{from} \StringTok{\textquotesingle{}@axona/protocol/persistence/file.js\textquotesingle{}}\OperatorTok{;}
\end{Highlighting}
\end{Shaded}

\subsubsection{15.9 Version constants}\label{version-constants}

\begin{Shaded}
\begin{Highlighting}[]
\NormalTok{WIRE\_VERSION         }\CommentTok{// \textquotesingle{}4.0\textquotesingle{}}
\NormalTok{KERNEL\_VERSION       }\CommentTok{// \textquotesingle{}4.11.2\textquotesingle{}}
\end{Highlighting}
\end{Shaded}

\subsubsection{15.10 Error codes worth
catching}\label{error-codes-worth-catching}

Errors subclass \texttt{AxonaError} with a stable \texttt{.code} --
switch on \texttt{.code}, not \texttt{.message}:

{\def\LTcaptype{none} % do not increment counter
\begin{longtable}[]{@{}
  >{\raggedright\arraybackslash}p{(\linewidth - 2\tabcolsep) * \real{0.5000}}
  >{\raggedright\arraybackslash}p{(\linewidth - 2\tabcolsep) * \real{0.5000}}@{}}
\toprule\noalign{}
\begin{minipage}[b]{\linewidth}\raggedright
Code
\end{minipage} & \begin{minipage}[b]{\linewidth}\raggedright
When
\end{minipage} \\
\midrule\noalign{}
\endhead
\bottomrule\noalign{}
\endlastfoot
\texttt{PUBLISH\_NO\_PUBLISH\_IDENTITY} & \texttt{pub} called without
\texttt{signWith} (13.1) \\
\texttt{PUBLISH\_INVALID\_TOPIC} & a bare Topic ID passed where a
descriptor is required (13.2) \\
\texttt{WRITE\_POLICY\_VIOLATION} & non-owner publishing to an owned
topic (13.5) \\
\texttt{PUBLISH\_PAYLOAD\_TOO\_LARGE} & enveloped message over the 16
KiB floor (13.7) \\
\texttt{PUBLISH\_INVALID\_MESSAGE} & message not JSON-serializable \\
\texttt{TOPIC\_REGION\_REQUIRED} & region omitted and no node region
available \\
\texttt{UPGRADE\_REQUIRED} & peer below the bridge's minimum wire
version (13.10) \\
\end{longtable}
}

\begin{center}\rule{0.5\linewidth}{0.5pt}\end{center}

\subsection{Where to go next}\label{where-to-go-next}

\begin{itemize}
\tightlist
\item
  \textbf{\href{Quick-Start-v4.11.2.md}{Quick Start}} -- a five-minute
  roundtrip for someone you are onboarding.
\item
  \textbf{\href{Axona-API-Reference-v4.11.2.md}{API Reference}} -- the
  exact signature of every public symbol.
\item
  \textbf{\href{Axona-Services-Guide-v0.3.md}{Services Guide}} -- the
  programs you \emph{run} rather than write: the signaling bridge, the
  relay + its front-ends (console, CLI, MCP server, desktop),
  directory/federation, and the PoW collector.
\item
  \textbf{\href{../architecture/Identity-and-Authorship-Model-v0.3.md}{Identity
  \& Authorship Model}} -- the design rationale behind the
  three-primitive model.
\item
  \textbf{Read \texttt{apps/axona-minimal/}} in the kernel repo -- a
  complete, runnable v3 app this guide's worked example is based on.
\item
  \textbf{Read the \href{../SECURITY-CHANGELOG.md}{security changelog}}
  -- the versioned record of what the protocol secures.
\end{itemize}

Found a bug or a gap in this guide? Open an issue at
\url{https://github.com/axona-net/axona-docs}.

\end{document}
